% MTL730 -- Lectures 2 through 8
\documentclass[11pt]{article}

\usepackage[a4paper,margin=1in]{geometry}
\usepackage{amsmath,amssymb,amsthm,mathtools}
\usepackage{enumitem}
\usepackage{hyperref}
\usepackage{microtype}

\hypersetup{
    colorlinks=true,
    linkcolor=blue,
    urlcolor=blue,
    citecolor=blue
}

\newtheorem{theorem}{Theorem}[section]
\newtheorem{lemma}[theorem]{Lemma}
\newtheorem{proposition}[theorem]{Proposition}
\newtheorem{corollary}[theorem]{Corollary}
\newtheorem{definition}[theorem]{Definition}
\newtheorem{example}[theorem]{Example}
\newtheorem{claim}[theorem]{Claim}
\theoremstyle{remark}
\newtheorem*{remark}{Remark}

\DeclareMathOperator{\gcdop}{gcd}
\DeclareMathOperator{\lcmop}{lcm}
\DeclareMathOperator{\ord}{ord}
\DeclareMathOperator{\ind}{ind}
\DeclareMathOperator{\degp}{deg}
\DeclareMathOperator{\charp}{char}
\newcommand{\Z}{\mathbb Z}
\newcommand{\N}{\mathbb N}
\newcommand{\Q}{\mathbb Q}
\newcommand{\R}{\mathbb R}
\newcommand{\F}{\mathbb F}
\newcommand{\Fp}{\mathbb F_p}
\newcommand{\phiE}{\varphi}
\newcommand{\divides}{\mid}
\newcommand{\notdivides}{\nmid}
\newcommand{\congmod}[1]{\pmod{#1}}
\newcommand{\Zmod}[1]{\mathbb Z/#1\mathbb Z}
\newcommand{\units}[1]{\left(\mathbb Z/#1\mathbb Z\right)^\times}

\title{\textbf{MTL730 --- Lectures 2--8}\\[2mm]
\large Number Theory and Finite Fields}
\author{}
\date{}

\begin{document}
\maketitle
\tableofcontents
\newpage
\section*{Lecture 2}
\addcontentsline{toc}{section}{Lecture 2}

% ============================================================
\section{Divisibility and Primes}
% ============================================================

\subsection{Divisibility}

\begin{definition}[Divisibility]
Let $a,b\in\Z$. We say that $a$ is a \emph{divisor} of $b$ if there exists
an integer $c$ such that
\[
    b=ac.
\]
If $a\neq 0$, the choice of $c$ is unique, and we write
\[
    c=\frac{b}{a}.
\]
We write $a\mid b$ when $a$ divides $b$, and $a\nmid b$ otherwise.
\end{definition}

\begin{remark}
The integer $c$ satisfying $b=ac$ is called the quotient of $b$ by $a$.
\end{remark}

\subsection{Common Divisors and the GCD}

\begin{definition}[Common divisor]
Let $a,b\in\Z$. An integer $d$ is a \emph{common divisor} of $a$ and $b$
if
\[
    d\mid a
    \qquad\text{and}\qquad
    d\mid b.
\]
\end{definition}

If $d$ is a common divisor of $a,b$, then for every $x,y\in\Z$,
\[
    d\mid ax+by.
\]
Indeed, if $a=da'$ and $b=db'$, then
\[
    ax+by=d(a'x+b'y).
\]

\begin{theorem}[B\'ezout-type characterization of the GCD]
For any two integers $a,b$, there exists a common divisor $d$ of $a,b$
such that
\[
    d=ax+by
\]
for some $x,y\in\Z$.
Moreover, every other common divisor of $a,b$ divides $d$.
\end{theorem}

\begin{proof}
We prove the existence by induction on $a+b$, after assuming without loss
of generality that
\[
    a\ge b\ge 0.
\]
If $a=b=0$, take $d=0$.

For the inductive step, if $b=0$, then $d=a$ works. Otherwise $b\ge1$.
Consider $a-b$. Since
\[
    (a-b)+b=a<a+b,
\]
the induction hypothesis applies to the pair $(a-b,b)$. Hence there is an
integer $d$ such that
\[
    d\mid(a-b),\qquad d\mid b,
\]
and
\[
    d=(a-b)x+by
\]
for some $x,y\in\Z$. Since
\[
    a=(a-b)+b,
\]
we also have $d\mid a$. Further,
\[
    d=ax+b(y-x),
\]
so $d$ is of the required form.

Finally, if $c$ is any common divisor of $a,b$, then
$c\mid ax+by=d$. Thus every common divisor divides $d$.
\end{proof}

\begin{corollary}
There is a unique non-negative common divisor $d$ of $a,b$ such that every
common divisor of $a,b$ divides $d$.
\end{corollary}

\begin{definition}[Greatest common divisor]
The unique non-negative integer in the preceding corollary is called the
\emph{greatest common divisor} of $a,b$, denoted by
\[
    \gcdop(a,b).
\]
If $(a,b)\neq(0,0)$, this is equivalently the least positive integer linear
combination of $a$ and $b$:
\[
    \gcdop(a,b)
    =
    \min\{\,d\in\Z : d>0,\ d=ax+by\text{ for some }x,y\in\Z\,\}.
\]
One sets $\gcdop(0,0)=0$ separately.
\end{definition}

\begin{remark}[Linear-combination characterization]
If $d=\gcdop(a,b)$, then there exist $x,y\in\Z$ such that
\[
    d=ax+by.
\]
Conversely, if $c=ax+by$ and $d=\gcdop(a,b)$, then
\[
    d\mid c.
\]
\end{remark}

\subsection{Relatively Prime Integers}

\begin{definition}[Relatively prime]
Two integers $a,b$ are said to be \emph{relatively prime} (or coprime) if
\[
    \gcdop(a,b)=1.
\]
\end{definition}

% ============================================================
\section{Division Algorithm}
% ============================================================

\begin{theorem}[Division Algorithm]
For any two given integers $a,b$ with $b>0$, there exists a unique pair
of integers $(q,r)$ such that
\[
    a=bq+r,
    \qquad
    0\le r<b.
\]
\end{theorem}

\begin{proof}[Proof sketch]
Consider the set
\[
    S=\{\,y\in\Z:y=a-bx,\ x\in\Z,\ y\ge0\,\}.
\]
The set $S$ is non-empty: choose $m\in\N$ sufficiently large that
\[
    mb>-a,
\]
so that
\[
    a+mb>0.
\]
Hence $S$ has a least element; call it $r$.

Then for some $q\in\Z$,
\[
    r=a-bq,
\]
or equivalently,
\[
    a=bq+r.
\]
It remains to show that $r<b$. If $r\ge b$, then
\[
    r-b=a-b(q+1)\ge0,
\]
so $r-b\in S$, contradicting the minimality of $r$. Thus
\[
    0\le r<b.
\]

For uniqueness, suppose
\[
    a=bq+r=bq'+r',
    \qquad
    0\le r,r'<b.
\]
Then
\[
    b(q-q')=r'-r.
\]
But
\[
    |r'-r|<b,
\]
so the only multiple of $b$ it can equal is $0$. Hence
\[
    r=r'
    \qquad\text{and}\qquad
    q=q'.
\]
\end{proof}

% ============================================================
\section{Euclidean Algorithm}
% ============================================================

\begin{lemma}
If
\[
    a=qb+r,
\]
then
\[
    \gcdop(a,b)=\gcdop(r,b).
\]
\end{lemma}

\begin{proof}
Let $d=\gcdop(a,b)$. Since $d\mid a$ and $d\mid b$,
\[
    d\mid(a-qb)=r.
\]
Thus $d$ is a common divisor of $r,b$.

Conversely, if $c\mid r$ and $c\mid b$, then
\[
    c\mid(qb+r)=a.
\]
Hence the common divisors of $(a,b)$ and $(r,b)$ are the same, so their
greatest common divisors are equal.
\end{proof}

\subsection{Finding the GCD}

Assume $a\ge b>0$. By the division algorithm,
\[
    a=q_1b+r_1,
    \qquad 0\le r_1<b.
\]
Then
\[
    \gcdop(a,b)=\gcdop(b,r_1).
\]
Continue:
\begin{align*}
    b &= q_2r_1+r_2,
        &0\le r_2<r_1,\\
    r_1 &= q_3r_2+r_3,
        &0\le r_3<r_2,\\
    &\ \vdots\\
    r_{n-2} &= q_nr_{n-1}+r_n,
        &0\le r_n<r_{n-1},\\
    r_{n-1} &= q_{n+1}r_n.
\end{align*}
The process terminates because the non-negative remainders form a strictly
decreasing sequence until the last non-zero remainder.

Using the lemma repeatedly,
\[
\begin{aligned}
    \gcdop(a,b)
    &=\gcdop(b,r_1)\\
    &=\gcdop(r_1,r_2)\\
    &=\cdots\\
    &=\gcdop(r_{n-1},r_n)\\
    &=r_n.
\end{aligned}
\]
Thus the last non-zero remainder is $\gcdop(a,b)$.

% ============================================================
\section{Prime Numbers}
% ============================================================

\begin{definition}[Prime number]
An integer $n>1$ is called \emph{prime} if its only positive divisors are
$1$ and $n$.

Equivalently, for every integer $a$,
\[
    a\mid n
    \quad\Longrightarrow\quad
    a=1\ \text{or}\ a=n
\]
when considering positive divisors.
\end{definition}

\subsection{Infinitely Many Primes}

\begin{theorem}[Euclid]
There are infinitely many prime numbers.
\end{theorem}

\begin{proof}
Suppose, for contradiction, that there are only finitely many primes,
say
\[
    p_1,p_2,\ldots,p_n.
\]
Consider
\[
    N=p_1p_2\cdots p_n+1.
\]
Let $p$ be any prime divisor of $N$. Then $p$ cannot be any of
$p_1,\ldots,p_n$, because for each $i$,
\[
    N\equiv1\pmod{p_i}.
\]
Thus $N$ has a prime divisor not in the assumed complete list, a
contradiction.
\end{proof}

\begin{remark}
The lecture also mentions the question of whether there are infinitely many
primes of the form
\[
    n^2+1.
\]
This is stated as a question in the lecture.
\end{remark}

\subsection{Fundamental Theorem of Arithmetic}

\begin{theorem}[Fundamental Theorem of Arithmetic]
Any integer $n>1$ can be written uniquely as a product of distinct prime
powers:
\[
    n=p_1^{e_1}p_2^{e_2}\cdots p_k^{e_k},
\]
where the $p_i$ are distinct primes and $e_i\ge1$.
\end{theorem}

% ============================================================
\section{Linear Diophantine Equations}
% ============================================================

\begin{theorem}[Solvability criterion]
For $a,b,c\in\Z$, the linear Diophantine equation
\[
    ax+by=c
\]
has a solution in $\Z^2$ if and only if
\[
    \gcdop(a,b)\mid c.
\]
\end{theorem}

\begin{proof}
Suppose first that
\[
    ax+by=c
\]
has an integer solution. Since $\gcdop(a,b)$ divides both $a$ and $b$,
it divides every integer linear combination of $a$ and $b$. Hence
\[
    \gcdop(a,b)\mid c.
\]

Conversely, suppose
\[
    \gcdop(a,b)\mid c.
\]
Let
\[
    d=\gcdop(a,b).
\]
By B\'ezout's identity, there exist $m,n\in\Z$ such that
\[
    d=am+bn.
\]
Since $d\mid c$, write
\[
    c=kd
\]
for some $k\in\Z$. Therefore
\[
    c=k(am+bn)=a(km)+b(kn),
\]
so
\[
    x=km,\qquad y=kn
\]
is an integer solution.
\end{proof}

\subsection{Example: $x+2y=5$}

Consider
\[
    x+2y=5.
\]
Since
\[
    \gcdop(1,2)=1,
\]
the equation has integer solutions. One solution is
\[
    (x,y)=(5,0).
\]

The corresponding homogeneous equation is
\[
    x+2y=0.
\]

If $(x_0,y_0)$ is one solution of $ax+by=c$, then every other solution
differs from it by a solution of the homogeneous equation
\[
    ax+by=0.
\]
Thus the solutions can be studied in a single framework.

% ============================================================
\section{Modular Arithmetic}
% ============================================================

\begin{definition}[Congruence modulo $m$]
Let $m\in\N$ and $a,b\in\Z$. We say that $a$ is \emph{congruent to $b$
modulo $m$}, denoted
\[
    a\equiv b\pmod m,
\]
if
\[
    m\mid(a-b).
\]
\end{definition}

\subsection{Properties of Congruence}

Let $m\in\N$ and $a,b,c,d\in\Z$.

\begin{enumerate}[label=\arabic*.]
    \item If
    \[
        a\equiv b\pmod m,
        \qquad
        c\equiv d\pmod m,
    \]
    then
    \[
        ac\equiv bd\pmod m.
    \]
    Also, for any $x,y\in\Z$,
    \[
        ax+cy\equiv bx+dy\pmod m.
    \]
    More generally, if $f$ is a polynomial with integer coefficients,
    then
    \[
        f(a)\equiv f(b)\pmod m.
    \]

    \item If
    \[
        a\equiv b\pmod m
    \]
    and $d>0$, then
    \[
        ad\equiv bd\pmod{md}.
    \]

    \item If
    \[
        ac\equiv bc\pmod m,
    \]
    and
    \[
        d=\gcdop(m,c),
    \]
    then
    \[
        a\equiv b\pmod{m/d}.
    \]
    In particular, if $\gcdop(m,c)=1$, then
    \[
        ac\equiv bc\pmod m
        \quad\Longrightarrow\quad
        a\equiv b\pmod m.
    \]

    \item If
    \[
        a\equiv b\pmod m,
    \]
    then
    \[
        \gcdop(a,m)=\gcdop(b,m).
    \]

    \item If
    \[
        a\equiv b\pmod m
    \]
    and
    \[
        0\le |a-b|<m,
    \]
    then
    \[
        a=b.
    \]
\end{enumerate}

\subsection{Congruence Classes}

\begin{proposition}
The congruence relation modulo $m$ is an equivalence relation on $\Z$:
\[
    a\,R\,b
    \quad\Longleftrightarrow\quad
    a\equiv b\pmod m.
\]
\end{proposition}

The equivalence class of $a$ is
\[
    \overline{a}
    =
    \{\,x\in\Z:x\equiv a\pmod m\,\}.
\]
This is called the \emph{residue class of $a$ modulo $m$}.

The residue classes
\[
    \overline{0},\overline{1},\overline{2},\ldots,\overline{m-1}
\]
are all distinct and disjoint, and
\[
    \Z
    =
    \overline{0}\cup\overline{1}\cup\cdots\cup\overline{m-1}.
\]

% ============================================================
\section{Complete and Reduced Residue Systems}
% ============================================================

\begin{definition}[Complete residue system]
A collection of representatives, one from each of the $m$ residue classes
modulo $m$, is called a \emph{complete residue system} (CRS) modulo $m$.

For example,
\[
    \{1,2,3,\ldots,m\}
\]
is a complete residue system modulo $m$.
\end{definition}

\begin{example}
If
\[
    \{a_1,\ldots,a_m\}
\]
is a complete residue system modulo $m$ and
\[
    \gcdop(k,m)=1,
\]
then
\[
    \{ka_1,\ldots,ka_m\}
\]
is also a complete residue system modulo $m$.
\end{example}

\subsection{Euler's Phi Function}

\begin{definition}[Euler's totient function]
For $m\in\N$, Euler's phi function is
\[
    \phiE(m)
    =
    \#\{\,a:1\le a\le m,\ \gcdop(a,m)=1\,\}.
\]
\end{definition}

\begin{remark}
The lecture writes the count using $1\le a\le m$; equivalently, for
$m>1$, one may use $1\le a<m$, since $m$ is not coprime to itself.
\end{remark}

\begin{definition}[Reduced residue system]
A \emph{reduced residue system} (RRS) modulo $m$ is a collection of
$\phiE(m)$ integers such that each integer is coprime to $m$ and no two
of them are congruent modulo $m$.
\end{definition}

\begin{example}
If
\[
    \{a_1,\ldots,a_{\phiE(m)}\}
\]
is an RRS modulo $m$ and
\[
    \gcdop(k,m)=1,
\]
then
\[
    \{ka_1,\ldots,ka_{\phiE(m)}\}
\]
is also an RRS modulo $m$.
\end{example}

% ============================================================
\section{Euler's Theorem}
% ============================================================

\begin{theorem}[Euler's theorem]
Let $m\in\N$ and let $a\in\Z$ satisfy
\[
    \gcdop(a,m)=1.
\]
Then
\[
    a^{\phiE(m)}\equiv1\pmod m.
\]
\end{theorem}

\begin{proof}
Let
\[
    \{a_1,a_2,\ldots,a_{\phiE(m)}\}
\]
be a reduced residue system modulo $m$.

For each $i$, multiplication by $a$ gives another residue coprime to $m$.
Moreover, the residues
\[
    aa_1,aa_2,\ldots,aa_{\phiE(m)}
\]
are pairwise incongruent modulo $m$.

Indeed, if for some $i,j$,
\[
    aa_i\equiv aa_j\pmod m,
\]
then, since $\gcdop(a,m)=1$, cancellation is valid and gives
\[
    a_i\equiv a_j\pmod m.
\]
Hence $i=j$.

Therefore
\[
    \{aa_1,\ldots,aa_{\phiE(m)}\}
\]
is also an RRS modulo $m$. Consequently,
\[
    aa_1\cdot aa_2\cdots aa_{\phiE(m)}
    \equiv
    a_1a_2\cdots a_{\phiE(m)}
    \pmod m.
\]
Thus
\[
    a^{\phiE(m)}
    a_1a_2\cdots a_{\phiE(m)}
    \equiv
    a_1a_2\cdots a_{\phiE(m)}
    \pmod m.
\]
Every $a_i$ is coprime to $m$, so the product is also coprime to $m$ and
may be cancelled. Hence
\[
    \boxed{a^{\phiE(m)}\equiv1\pmod m}.
\]
\end{proof}

% ============================================================
\section{Fermat's Little Theorem}
% ============================================================

\begin{theorem}[Fermat's Little Theorem]
Let $p$ be a prime number and let $a\in\Z$. Then
\[
    a^p\equiv a\pmod p.
\]
\end{theorem}

\begin{proof}
If $p\mid a$, then
\[
    a^p\equiv a\equiv0\pmod p.
\]

If $p\not\mid a$, then
\[
    \gcdop(a,p)=1.
\]
By Euler's theorem,
\[
    a^{\phiE(p)}\equiv1\pmod p.
\]
Since $p$ is prime,
\[
    \phiE(p)=p-1.
\]
Thus
\[
    a^{p-1}\equiv1\pmod p.
\]
Multiplying by $a$ gives
\[
    a^p\equiv a\pmod p.
\]
\end{proof}

% ============================================================
\section{Computing Euler's Phi Function}
% ============================================================

For
\[
    m=p_1^{e_1}p_2^{e_2}\cdots p_r^{e_r},
\]
where the $p_i$ are distinct primes, the lecture derives
\[
    \boxed{
    \phiE(m)
    =
    m\prod_{p\mid m}\left(1-\frac1p\right)
    }.
\]

\subsection{Two Claims}

\begin{claim}
If
\[
    \gcdop(m,n)=1,
\]
then
\[
    \phiE(mn)=\phiE(m)\phiE(n).
\]
\end{claim}

\begin{claim}
If $p$ is prime and $e\ge1$, then
\[
    \phiE(p^e)=p^e-p^{e-1}
    =
    p^e\left(1-\frac1p\right).
\]
\end{claim}

\subsection{Proof of Claim 2}

Among
\[
    \{1,2,\ldots,p^e\},
\]
an integer is not coprime to $p^e$ precisely when it is divisible by $p$.
There are
\[
    p^{e-1}
\]
multiples of $p$ in this set. Therefore
\[
    \phiE(p^e)
    =
    p^e-p^{e-1}.
\]

\subsection{Proof of Claim 1}

\begin{proof}
Assume
\[
    \gcdop(m,n)=1.
\]
Consider the $n\times m$ array whose rows are indexed by $0\le k\le n-1$
and whose columns are indexed by $1\le\ell\le m$:
\[
\begin{pmatrix}
1 & 2 & 3 & \cdots & m\\
m+1 & m+2 & m+3 & \cdots & 2m\\
2m+1 & 2m+2 & 2m+3 & \cdots & 3m\\
\vdots & \vdots & \vdots & \ddots & \vdots\\
(n-1)m+1 & (n-1)m+2 & (n-1)m+3 & \cdots & nm
\end{pmatrix}
\]
with entries
\[
    km+\ell,
    \qquad
    0\le k\le n-1,\quad 1\le\ell\le m.
\]

The entries consist of a complete set of integers from $1$ through $mn$.

We want to count the entries coprime to $mn$. Since
\[
    \gcdop(m,n)=1,
\]
an integer $x$ is coprime to $mn$ exactly when
\[
    \gcdop(x,m)=1
    \qquad\text{and}\qquad
    \gcdop(x,n)=1.
\]

Fix a column $\ell$ that is coprime to $m$. If two entries in this column,
\[
    k_1m+\ell
    \qquad\text{and}\qquad
    k_2m+\ell,
\]
are congruent modulo $n$, then
\[
    k_1m+\ell\equiv k_2m+\ell\pmod n,
\]
so
\[
    (k_1-k_2)m\equiv0\pmod n.
\]
Because $\gcdop(m,n)=1$,
\[
    k_1-k_2\equiv0\pmod n.
\]
But
\[
    0\le k_1,k_2\le n-1,
\]
hence
\[
    k_1=k_2.
\]
Thus the $n$ entries in such a column give distinct residue classes modulo $n$.

Exactly $\phiE(n)$ of them are coprime to $n$. There are
$\phiE(m)$ columns whose $\ell$ is coprime to $m$. Hence the total number
of entries coprime to both $m$ and $n$ is
\[
    \phiE(m)\phiE(n).
\]
Therefore
\[
    \boxed{\phiE(mn)=\phiE(m)\phiE(n)}.
\]
\end{proof}

\subsection{Deriving the Product Formula}

Write the prime factorization of $m$ as
\[
    m=p_1^{e_1}p_2^{e_2}\cdots p_r^{e_r}.
\]
Since the prime powers are pairwise coprime, repeated application of
Claim 1 gives
\[
    \phiE(m)
    =
    \phiE(p_1^{e_1})
    \phiE(p_2^{e_2})
    \cdots
    \phiE(p_r^{e_r}).
\]
Using Claim 2,
\begin{align*}
    \phiE(m)
    &=
    \prod_{i=1}^r
    \left(p_i^{e_i}-p_i^{e_i-1}\right)\\
    &=
    \prod_{i=1}^r
    p_i^{e_i}\left(1-\frac1{p_i}\right)\\
    &=
    \left(\prod_{i=1}^r p_i^{e_i}\right)
    \prod_{i=1}^r
    \left(1-\frac1{p_i}\right)\\
    &=
    m\prod_{p\mid m}
    \left(1-\frac1p\right).
\end{align*}

Thus
\[
    \boxed{
    \phiE(m)
    =
    m\prod_{p\mid m}
    \left(1-\frac1p\right)
    }.
\]

% ============================================================
\section*{Summary of the Lecture}
% ============================================================

The lecture develops the following chain of ideas:
\[
\boxed{
\text{Divisibility}
\longrightarrow
\gcd
\longrightarrow
\text{Division Algorithm}
\longrightarrow
\text{Euclidean Algorithm}
\longrightarrow
\text{Primes}
}
\]
followed by
\[
\boxed{
\text{Diophantine Equations}
\longrightarrow
\text{Modular Arithmetic}
\longrightarrow
\text{Residue Systems}
\longrightarrow
\phiE(m)
\longrightarrow
\text{Euler's Theorem}
\longrightarrow
\text{Fermat's Little Theorem}
}.
\]

The key formulas are
\[
    \gcdop(a,b)=\gcdop(b,r)
    \quad\text{when }a=qb+r,
\]
\[
    ax+by=c
    \quad\Longleftrightarrow\quad
    \gcdop(a,b)\mid c,
\]
\[
    a^{\phiE(m)}\equiv1\pmod m
    \qquad
    \left(\gcdop(a,m)=1\right),
\]
\[
    a^p\equiv a\pmod p
    \qquad
    (p\text{ prime}),
\]
and
\[
    \phiE(m)
    =
    m\prod_{p\mid m}\left(1-\frac1p\right).
\]

\clearpage

\section*{Lecture 3 (continued)}
\addcontentsline{toc}{section}{Lecture 3 (continued)}

% ============================================================
\section{Linear Congruences}
% ============================================================

We first recall the familiar case
\[
    ax\equiv b\pmod m,
\]
where
\[
    \gcdop(a,m)=1.
\]
Since $a$ is invertible modulo $m$,
\[
    x\equiv a^{\phiE(m)-1}b\pmod m.
\]
Thus there is a unique solution modulo $m$.

\subsection{The General Case}

The natural question is:

\begin{center}
\emph{What if $\gcdop(a,m)\neq1$?}
\end{center}

Consider
\[
    ax\equiv b\pmod m.
\]
This is equivalent to
\[
    ax+my=b.
\]

Let
\[
    d=\gcdop(a,m).
\]
Since
\[
    d\mid a,\qquad d\mid m,
\]
every integer of the form
\[
    ax+my
\]
is divisible by $d$. Hence a necessary condition for solvability is
\[
    d\mid b.
\]

\begin{theorem}
The linear congruence
\[
    \boxed{ax\equiv b\pmod m}
\]
has a solution if and only if
\[
    \boxed{\gcdop(a,m)\mid b.}
\]
If
\[
    d=\gcdop(a,m)\mid b,
\]
then the congruence has exactly $d$ distinct solutions modulo $m$.
\end{theorem}

\subsection{Description of the Solutions}

Suppose
\[
    d=\gcdop(a,m)\mid b.
\]
Then the congruence can be divided by $d$:
\[
    \frac ad x
    \equiv
    \frac bd
    \pmod{\frac md}.
\]
Since
\[
    \gcdop\!\left(\frac ad,\frac md\right)=1,
\]
there is a unique solution modulo $m/d$. Let this solution be $t$:
\[
    \frac ad\,t
    \equiv
    \frac bd
    \pmod{\frac md}.
\]

The complete set of solutions modulo $m$ is then
\[
    \boxed{
    x\equiv t+\frac{km}{d}\pmod m,
    \qquad
    k=0,1,\ldots,d-1.
    }
\]

\subsection{Proof That the Solutions Are Distinct}

Suppose
\[
    t+\frac{km}{d}
    \equiv
    t+\frac{\ell m}{d}
    \pmod m.
\]
Then
\[
    \frac{(k-\ell)m}{d}\equiv0\pmod m.
\]
Equivalently,
\[
    m\mid\frac{(k-\ell)m}{d},
\]
so
\[
    d\mid(k-\ell).
\]
But
\[
    0\le k,\ell\le d-1,
\]
hence
\[
    k=\ell.
\]
Therefore the $d$ displayed solutions are pairwise distinct modulo $m$.

\subsection{Example}

Consider
\[
    5x\equiv3\pmod{24}.
\]
Since
\[
    \gcdop(5,24)=1,
\]
we may use the inverse of $5$ modulo $24$:
\[
    5^{-1}\equiv5\pmod{24}.
\]
Therefore
\[
    x\equiv5\cdot3
    \equiv15\pmod{24}.
\]

Thus the solution is
\[
    \boxed{x\equiv15\pmod{24}.}
\]

As another example, consider
\[
    25x\equiv15\pmod{120}.
\]
Here
\[
    \gcdop(25,120)=5
\]
and
\[
    5\mid15.
\]
Hence there are exactly $5$ solutions modulo $120$.

Dividing the congruence by $5$ gives
\[
    5x\equiv3\pmod{24},
\]
so
\[
    x\equiv15\pmod{24}.
\]
Therefore the five solutions modulo $120$ are
\[
    \boxed{
    15,\ 39,\ 63,\ 87,\ 111.
    }
\]

% ============================================================
\section{Systems of Linear Congruences}
% ============================================================

A system of linear congruences is a family such as
\[
    \begin{cases}
        x\equiv b_1\pmod{m_1},\\
        x\equiv b_2\pmod{m_2},\\
        \vdots\\
        x\equiv b_k\pmod{m_k}.
    \end{cases}
\]

\begin{theorem}[Chinese Remainder Theorem]
Let
\[
    m_1,m_2,\ldots,m_k
\]
be positive integers such that every two of them are coprime:
\[
    \gcdop(m_i,m_j)=1
    \qquad(i\neq j).
\]
Then for any given integers
\[
    b_1,b_2,\ldots,b_k,
\]
the system
\[
    \boxed{
    x\equiv b_i\pmod{m_i},
    \qquad
    1\le i\le k
    }
\]
has a unique solution modulo
\[
    M=m_1m_2\cdots m_k.
    \]
\end{theorem}

\subsection{Uniqueness}

Suppose $x$ and $y$ are two solutions. Then for every $i$,
\[
    x\equiv y\pmod{m_i}.
\]
Thus
\[
    m_i\mid(x-y)
    \qquad\text{for every }i.
\]
Since the moduli are pairwise coprime,
\[
    M=m_1m_2\cdots m_k
\]
also divides $x-y$. Hence
\[
    x\equiv y\pmod M.
\]
Therefore the solution is unique modulo $M$.

\subsection{Existence}

Let
\[
    M=m_1m_2\cdots m_k,
    \qquad
    M_i=\frac{M}{m_i}.
\]
Since $m_i$ is coprime to every other modulus,
\[
    \gcdop(M_i,m_i)=1.
\]
Hence there exists an integer $N_i$ such that
\[
    \boxed{
    M_iN_i\equiv1\pmod{m_i}.
    }
\]
Equivalently, there exist integers $N_i,r_i$ satisfying
\[
    \boxed{
    1=M_iN_i+m_ir_i.
    }
\]

Define
\[
    \boxed{
    x=\sum_{i=1}^k b_iM_iN_i.
    }
\]

Fix $j$. Modulo $m_j$, every term with $i\neq j$ vanishes because
\[
    m_j\mid M_i.
\]
Therefore
\[
    x
    \equiv
    b_jM_jN_j
    \equiv
    b_j
    \pmod{m_j}.
\]
Thus
\[
    x\equiv b_j\pmod{m_j}
\]
for every $j$, proving existence.

Hence
\[
    \boxed{
    x\equiv
    \sum_{i=1}^k b_iM_iN_i
    \pmod M
    }.
\]

\subsection{Constructing the $N_i$}

The integers $N_i$ can be found using the Euclidean algorithm.

Since
\[
    \gcdop(M_i,m_i)=1,
\]
Bezout's identity gives
\[
    M_iN_i+m_ir_i=1.
\]
Reducing modulo $m_i$ gives
\[
    M_iN_i\equiv1\pmod{m_i}.
\]

% ============================================================
\section{Polynomial Congruences}
% ============================================================

The Chinese Remainder Theorem also gives a useful result for polynomial
congruences.

Let
\[
    m_1,m_2,\ldots,m_k
\]
be pairwise coprime positive integers, and let
\[
    M=m_1m_2\cdots m_k.
\]

\begin{corollary}
Let $f(x)\in\Z[x]$. Then
\[
    f(x)\equiv0\pmod M
\]
has a solution if and only if
\[
    f(x)\equiv0\pmod{m_i}
\]
has a solution for every $i$.
\end{corollary}

Indeed, if
\[
    x_i
\]
is a solution modulo $m_i$, then the Chinese Remainder Theorem gives a
simultaneous solution
\[
    x\equiv x_i\pmod{m_i}
    \qquad(1\le i\le k).
\]
Since polynomial congruences respect congruence,
\[
    f(x)\equiv f(x_i)\equiv0\pmod{m_i}
\]
for every $i$. Hence
\[
    f(x)\equiv0\pmod M.
\]

Moreover, if
\[
    f(x)\equiv0\pmod{m_i}
\]
has
\[
    s(m_i)
\]
solutions modulo $m_i$, then
\[
    \boxed{
    f(x)\equiv0\pmod M
    }
\]
has
\[
    \boxed{
    \prod_{i=1}^k s(m_i)
    }
\]
solutions modulo $M$.

% ============================================================
\section{A Bound on the Number of Roots of a Polynomial Modulo a Prime}
% ============================================================

Consider
\[
    x^2\equiv1\pmod3.
\]
The solutions are
\[
    x\equiv1,2\pmod3.
\]

On the other hand,
\[
    x^2\equiv1\pmod8
\]
has four solutions:
\[
    \boxed{x\equiv1,3,5,7\pmod8.}
\]

Thus the familiar statement that a polynomial of degree $n$ has at most
$n$ roots is not automatically true modulo an arbitrary composite
modulus.

The prime-modulus case behaves differently.

\begin{theorem}[Lagrange]
Let $p$ be a prime and let
\[
    f(x)
    =
    a_0+a_1x+\cdots+a_nx^n
    \in\Z[x],
\]
with
\[
    a_n\not\equiv0\pmod p.
\]
Then
\[
    f(x)\equiv0\pmod p
\]
has at most $n$ solutions modulo $p$.
\end{theorem}

\begin{proof}
We prove the result by induction on $n$.

For
\[
    n=0,
\]
we have
\[
    f(x)\equiv a_0\not\equiv0\pmod p,
\]
so there are no solutions.

Now assume the statement holds for degree $n-1$ polynomials.

If
\[
    f(x)\equiv0\pmod p
\]
has no solution, we are done. Suppose it has a solution $x_0$, and let $r$
be the total number of solutions modulo $p$. For any $x$,
\begin{align*}
    f(x)-f(x_0)
    &=
    \sum_{i=1}^n a_i(x^i-x_0^i)\\
    &=
    (x-x_0)g(x),
\end{align*}
where $g(x)$ is a polynomial of degree at most $n-1$.

Let
\[
    x_1,\ldots,x_{r-1}
\]
be the remaining solutions of
\[
    f(x)\equiv0\pmod p
\]
(so there are $r-1$ of them if $r\ge 1$). For each such $x_i$,
\[
    f(x_i)-f(x_0)\equiv0\pmod p.
\]
Therefore
\[
    (x_i-x_0)g(x_i)\equiv0\pmod p.
\]
Since
\[
    x_i\not\equiv x_0\pmod p,
\]
we must have
\[
    g(x_i)\equiv0\pmod p.
\]
Thus $g$ has at least $r-1$ roots. By the induction hypothesis,
\[
    r-1\le n-1.
\]
Hence
\[
    r\le n.
\]
Therefore $f$ has at most $n$ roots modulo $p$.
\end{proof}

% ============================================================
\section{A Useful Application}
% ============================================================

Let $p$ be a prime. Consider
\[
    f(x)
    =
    (x-1)(x-2)\cdots(x-(p-1)).
\]
Modulo $p$, this polynomial has the $p-1$ distinct roots
\[
    1,2,\ldots,p-1.
\]

Since
\[
    \deg f=p-1,
\]
Lagrange's theorem implies that these are all of its roots modulo $p$.
The two monic polynomials $f(x)$ and $x^{p-1}-1$ of degree $p-1$ therefore
have the same $p-1$ roots modulo $p$. Their difference has degree at most
$p-2$ but has $p-1$ roots, so it is the zero polynomial modulo $p$. Thus
\[
    \boxed{
    (x-1)(x-2)\cdots(x-(p-1))
    \equiv
    x^{p-1}-1
    \pmod p.
    }
\]

Comparing coefficients, all intermediate coefficients are divisible by
$p$:
\[
    \boxed{
    a_i\equiv0\pmod p,
    \qquad
    0<i<p-1.
    }
\]

The constant term is also especially useful.

% ============================================================
\section{Wilson's Theorem}
% ============================================================

\begin{theorem}[Wilson]
Let $p$ be a prime. Then
\[
    \boxed{
    (p-1)!\equiv-1\pmod p.
    }
\]
\end{theorem}

\begin{proof}
The case
\[
    p=2
\]
is immediate:
\[
    1!\equiv-1\pmod2.
\]

Now let $p$ be an odd prime. From the polynomial identity above,
\[
    (x-1)(x-2)\cdots(x-(p-1))
    \equiv
    x^{p-1}-1
    \pmod p.
\]

The constant term on the left is
\[
    (-1)(-2)\cdots(-(p-1)).
\]
Since $p-1$ is even,
\[
    (-1)^{p-1}(p-1)!
    =
    (p-1)!.
\]
The constant term on the right is
\[
    -1.
\]
Therefore
\[
    (p-1)!\equiv-1\pmod p.
\]
\end{proof}

\subsection{The Classical Pairing Proof}

There is also a direct proof. For every
\[
    a\in\{1,2,\ldots,p-1\},
\]
there is a unique inverse
\[
    a^{-1}\pmod p.
\]
The elements satisfying
\[
    a\equiv a^{-1}\pmod p
\]
are precisely the solutions of
\[
    a^2\equiv1\pmod p,
\]
namely
\[
    a\equiv\pm1\pmod p.
\]

All other elements can therefore be paired with their distinct inverses.
Their products contribute $1$ modulo $p$. Hence
\[
    (p-1)!
    \equiv
    1\cdot(-1)
    \equiv-1
    \pmod p.
\]

% ============================================================
\section{Converse of Wilson's Theorem}
% ============================================================

\begin{theorem}[Converse]
If $n>1$ and
\[
    \boxed{
    (n-1)!\equiv-1\pmod n,
    }
\]
then $n$ is prime.
\end{theorem}

\begin{proof}
Suppose, for contradiction, that $n$ is composite.

Then there exists an integer $d$ such that
\[
    1<d<n
\]
and
\[
    d\mid n.
\]
But
\[
    d\le n-1,
\]
so $d$ occurs among the factors of
\[
    (n-1)!.
\]
Hence
\[
    d\mid(n-1)!.
\]
On the other hand,
\[
    (n-1)!\equiv-1\pmod n
\]
implies
\[
    \gcdop((n-1)!,n)=1.
\]
This contradicts $d>1$ dividing both $(n-1)!$ and $n$.

Therefore $n$ cannot be composite, and hence
\[
    \boxed{n\text{ is prime}.}
\]
\end{proof}

% ============================================================
\section*{Summary}
% ============================================================

The lecture develops the following progression:
\[
\boxed{
\text{Linear congruences}
\longrightarrow
\text{CRT}
\longrightarrow
\text{Polynomial congruences}
\longrightarrow
\text{Lagrange's theorem}
\longrightarrow
\text{Wilson's theorem}
}.
\]

The main results are

\[
    ax\equiv b\pmod m
    \text{ is solvable}
    \quad\Longleftrightarrow\quad
    \gcdop(a,m)\mid b,
\]
and if
\[
    d=\gcdop(a,m)\mid b,
\]
then there are exactly $d$ solutions modulo $m$:
\[
    \boxed{
    x\equiv t+\frac{km}{d}\pmod m,
    \qquad
    k=0,\ldots,d-1.
    }
\]

For pairwise coprime moduli,
\[
    M=m_1\cdots m_k,
\]
the Chinese Remainder Theorem gives
\[
    \boxed{
    x\equiv
    \sum_{i=1}^k b_iM_iN_i
    \pmod M,
    \qquad
    M_i=\frac{M}{m_i},
    \quad
    M_iN_i\equiv1\pmod{m_i}.
    }
\]

For a polynomial of degree $n$ over $\Fp$ with nonzero leading
coefficient,
\[
    f(x)\equiv0\pmod p
\]
has at most $n$ solutions.

Finally,
\[
    \boxed{
    p\text{ is prime}
    \quad\Longleftrightarrow\quad
    (p-1)!\equiv-1\pmod p.
    }
\]

\clearpage

\section*{Lecture 4}
\addcontentsline{toc}{section}{Lecture 4}

% ============================================================
\section{Application of Wilson's Theorem}
% ============================================================

\begin{theorem}
Let $p$ be an odd prime. The congruence
\[
    x^2\equiv -1\pmod p
\]
has a solution if and only if
\[
    p\equiv 1\pmod 4.
\]
\end{theorem}

\begin{proof}
First suppose there is a solution: $a^2\equiv -1\pmod p$. Then $p\nmid a$, so
Fermat's little theorem gives
\[
    a^{p-1}\equiv 1\pmod p.
\]
Hence
\[
    \bigl(a^2\bigr)^{(p-1)/2}\equiv 1\pmod p,
\]
and therefore
\[
    (-1)^{(p-1)/2}\equiv 1\pmod p.
\]
If $p\equiv 3\pmod 4$, write $p=4n+3$. Then
\[
    \frac{p-1}{2}=2n+1
\]
is odd, so $(-1)^{(p-1)/2}=-1$, a contradiction. Thus $p\equiv 1\pmod 4$.

Conversely, suppose $p\equiv 1\pmod 4$, so $p=4n+1$. Wilson's theorem yields
\[
    (p-1)!\equiv -1\pmod p,
\]
that is
\[
    (4n)!\equiv -1\pmod p.
\]
Pair the factors $k$ and $p-k$:
\[
    p-k\equiv -k\pmod p.
\]
There are $2n$ such pairs, and $2n$ is even, so
\[
    1\cdot 2\cdots (2n)\cdot(2n+1)\cdots(4n)
    \equiv
    (2n)!\,(-1)^{2n}(2n)!
    \equiv
    \bigl((2n)!\bigr)^2
    \pmod p.
\]
Hence
\[
    \bigl((2n)!\bigr)^2\equiv -1\pmod p.
\]
Taking $x=(2n)!$ gives a solution of $x^2\equiv -1\pmod p$.
\end{proof}

% ============================================================
\section{Quadratic Congruences}
% ============================================================

Let $p$ be an odd prime. We wish to study
\[
    x^2\equiv n\pmod p
\]
and, more generally, the quadratic congruence
\[
    ax^2+bx+c\equiv 0\pmod p,
    \qquad p\nmid a.
\]
Multiplying by $4a$ (which is invertible modulo $p$) gives the equivalent
congruence
\[
    4a^2x^2+4abx+4ac\equiv 0\pmod p,
\]
or
\[
    (2ax+b)^2\equiv b^2-4ac\pmod p.
\]
Setting $Y=2ax+b$ and $n=b^2-4ac$, we are reduced to
\[
    Y^2\equiv n\pmod p.
\]

\subsection{Two solutions when one exists}

Assume $p$ is odd and $n\not\equiv 0\pmod p$. If $a$ is a solution of
$x^2\equiv n\pmod p$, then so is $p-a$, because
\[
    (p-a)^2=p^2-2ap+a^2\equiv a^2\pmod p.
\]
These two solutions are distinct: $a\equiv p-a\pmod p$ would force
$p\mid 2a$, hence $p\mid a$ (since $p$ is odd), contradicting $n\not\equiv 0$.

% ============================================================
\section{Quadratic Residues}
% ============================================================

\begin{definition}[Quadratic residue]
Let $n\in\Z$ with $n\not\equiv 0\pmod p$. We say that $n$ is a
\emph{quadratic residue} modulo $p$ if
\[
    x^2\equiv n\pmod p
\]
has a solution. Otherwise $n$ is a \emph{quadratic non-residue} modulo $p$.
\end{definition}

\begin{example}
For $p=11$,
\[
\begin{array}{c|cccccccccc}
a&1&2&3&4&5&6&7&8&9&10\\
\hline
a^2&1&4&9&5&3&3&5&9&4&1
\end{array}
\pmod{11}.
\]
The quadratic residues modulo $11$ are
\[
    \{1,3,4,5,9\},
\]
and the non-residues are
\[
    \{2,6,7,8,10\}.
\]
\end{example}

Since $a^2\equiv(p-a)^2\pmod p$, the distinct nonzero squares modulo $p$
are represented by
\[
    1^2,2^2,\ldots,\Bigl(\frac{p-1}{2}\Bigr)^2.
\]
Thus there are at most $(p-1)/2$ quadratic residues.

\begin{theorem}
Let $p$ be an odd prime. Among $1,2,\ldots,p-1$ there are exactly
$(p-1)/2$ quadratic residues and $(p-1)/2$ quadratic non-residues modulo
$p$. The residues are precisely
\[
    1^2,2^2,\ldots,\Bigl(\frac{p-1}{2}\Bigr)^2\pmod p.
\]
\end{theorem}

\begin{proof}
The displayed squares are pairwise incongruent. Indeed, if
\[
    a^2\equiv b^2\pmod p
    \qquad\text{with}\qquad
    1\le a,b\le\frac{p-1}{2},
\]
then
\[
    (a-b)(a+b)\equiv 0\pmod p,
\]
so $p\mid(a-b)$ or $p\mid(a+b)$. The bounds force $a=b$. Hence there are
exactly $(p-1)/2$ distinct nonzero squares, and therefore equally many
non-residues.
\end{proof}

% ============================================================
\section{The Legendre Symbol}
% ============================================================

\begin{definition}[Legendre symbol]
Let $p$ be an odd prime. For any integer $n$, set
\[
\left(\frac np\right)
=
\begin{cases}
0,& n\equiv 0\pmod p,\\[1mm]
1,& n\text{ is a quadratic residue modulo }p,\\[1mm]
-1,& n\text{ is a quadratic non-residue modulo }p.
\end{cases}
\]
For $p=2$ the lecture uses the convention
\[
\left(\frac n2\right)
=
\begin{cases}
0,& n\text{ even},\\[1mm]
1,& n\text{ odd}.
\end{cases}
\]
\end{definition}

The value $\bigl(\frac np\bigr)$ depends only on $n\bmod p$.

% ============================================================
\section{Euler's Criterion}
% ============================================================

\begin{theorem}[Euler's criterion]
Let $p$ be an odd prime. Then
\[
    \boxed{
    \left(\frac np\right)
    \equiv
    n^{(p-1)/2}
    \pmod p.
    }
\]
\end{theorem}

\begin{proof}
If $p\mid n$, both sides are $0$ modulo $p$.

Now assume $p\nmid n$. By Fermat, $n^{p-1}\equiv 1\pmod p$, so
\[
    \bigl(n^{(p-1)/2}\bigr)^2\equiv 1\pmod p.
\]
The congruence $X^2\equiv 1\pmod p$ has only the two solutions $X\equiv\pm 1$,
hence
\[
    n^{(p-1)/2}\equiv\pm 1\pmod p.
\]

If $\bigl(\frac np\bigr)=1$, then $n\equiv a^2\pmod p$ for some $a$ with
$p\nmid a$. Thus
\[
    n^{(p-1)/2}\equiv a^{p-1}\equiv 1\pmod p.
\]
There are exactly $(p-1)/2$ quadratic residues, and each is a root of
\[
    X^{(p-1)/2}-1\equiv 0\pmod p.
\]
A polynomial of degree $(p-1)/2$ over the field $\Z/p\Z$ has at most
$(p-1)/2$ roots, so these residues are \emph{all} the roots. Therefore if
$\bigl(\frac np\bigr)=-1$, we cannot have $n^{(p-1)/2}\equiv 1$, and must have
\[
    n^{(p-1)/2}\equiv -1\pmod p.
\]
\end{proof}

\begin{corollary}[Complete multiplicativity]
The map $n\mapsto\bigl(\frac np\bigr)$ is completely multiplicative:
\[
    \left(\frac{mn}{p}\right)
    =
    \left(\frac mp\right)
    \left(\frac np\right).
\]
\end{corollary}

\begin{proof}
Both sides are congruent to $(mn)^{(p-1)/2}=m^{(p-1)/2}n^{(p-1)/2}$
modulo $p$, and both sides take values in $\{0,\pm 1\}$.
\end{proof}

\begin{corollary}
\[
\left(\frac{-1}{p}\right)
=
(-1)^{(p-1)/2}
=
\begin{cases}
1,& p\equiv 1\pmod 4,\\[1mm]
-1,& p\equiv 3\pmod 4.
\end{cases}
\]
\end{corollary}

This recovers the solvability criterion for $x^2\equiv -1\pmod p$.

\begin{remark}[Two classical facts]
\begin{enumerate}[label=\arabic*.]
    \item An odd prime $p$ is a sum of two squares, $p=x^2+y^2$, if and only if
    $p\equiv 1\pmod 4$. (For instance $2=1^2+1^2$.)

    \item The only rational primes that remain prime in the Gaussian integers
    $\Z[i]$ are the primes $p\equiv 3\pmod 4$.
\end{enumerate}
\end{remark}

% ============================================================
\section{The Supplementary Law for $2$}
% ============================================================

\begin{theorem}
Let $p$ be an odd prime. Then
\[
\left(\frac 2p\right)
=
(-1)^{(p^2-1)/8}
=
\begin{cases}
1,& p\equiv\pm 1\pmod 8,\\[1mm]
-1,& p\equiv\pm 3\pmod 8.
\end{cases}
\]
\end{theorem}

\begin{proof}
Write $N=(p-1)/2$. Among the even residues $2,4,\ldots,2N=p-1$, replace each
term $2k>p/2$ by $2k-p\equiv -(p-2k)\pmod p$. The integers $p-2k$ arising
in this way, together with the even numbers that were already less than
$p/2$, are a rearrangement of $\{1,2,\ldots,N\}$. If $m$ is the number of
replacements, then
\[
    2\cdot 4\cdots(p-1)
    \equiv
    (-1)^m N!
    \pmod p.
\]
The left-hand side is $2^N N!$, so
\[
    2^N\equiv(-1)^m\pmod p.
\]
The integers $k\in\{1,\ldots,N\}$ with $2k>p/2$ are those with $k>p/4$, and
a check of the four odd classes modulo $8$ gives
\[
    m\equiv\frac{p^2-1}{8}\pmod 2.
\]
Euler's criterion now yields
\[
    \left(\frac 2p\right)
    \equiv
    2^{(p-1)/2}
    \equiv
    (-1)^{(p^2-1)/8}
    \pmod p.
\]
(The same count of $m$ is the special case $n=2$ of Gauss's lemma below.)
\end{proof}

\begin{corollary}
\[
\left(\frac{-2}{p}\right)
=
\begin{cases}
1,& p\equiv 1,3\pmod 8,\\[1mm]
-1,& p\equiv 5,7\pmod 8.
\end{cases}
\]
\end{corollary}

\begin{proof}
Multiply the formulae for $\bigl(\frac{-1}{p}\bigr)$ and $\bigl(\frac 2p\bigr)$.
\end{proof}

\begin{remark}
The lecture records that an odd prime $p$ has the form $p=x^2+2y^2$ if and
only if $p\equiv 1,3\pmod 8$. This is consistent with the evaluation of
$\bigl(\frac{-2}{p}\bigr)$, but the representation theorem itself is not
proved here. A standard reference is Cox, \emph{Primes of the Form
$x^2+ny^2$}.
\end{remark}

% ============================================================
\section{Gauss's Lemma}
% ============================================================

\begin{theorem}[Gauss's lemma]
Let $p$ be an odd prime and $p\nmid n$. Consider the least positive residues
modulo $p$ of
\[
    n,\ 2n,\ \ldots,\ \frac{p-1}{2}n.
\]
If exactly $m$ of these residues are greater than $p/2$, then
\[
    \boxed{\left(\frac np\right)=(-1)^m.}
\]
\end{theorem}

\begin{example}
The integer $4$ is a quadratic residue modulo $11$, and $\bigl(\frac{4}{11}\bigr)=1$.
The residues of $4,8,12,16,20$ modulo $11$ are
\[
    4,8,1,5,9.
\]
Exactly two of them ($8$ and $9$) exceed $11/2$, so $m=2$ and
$(-1)^m=1$.
\end{example}

\begin{proof}
Let $r_1,\ldots,r_{(p-1)/2}$ be the least positive residues of
$n,2n,\ldots,\frac{p-1}{2}n$ modulo $p$. These residues are pairwise distinct:
if $in\equiv jn\pmod p$ with $1\le i<j\le(p-1)/2$, then $p\mid n(j-i)$, hence
$p\mid(j-i)$, which is impossible.

Split them into two sets:
\[
    \{a_1,\ldots,a_\ell\}
    =
    \{r_k:r_k<p/2\},
    \qquad
    \{b_1,\ldots,b_m\}
    =
    \{r_k:r_k>p/2\},
\]
so $\ell+m=(p-1)/2$. Put $c_j=p-b_j$. Then $0<c_j<p/2$.

\begin{claim}
The integers $a_i$ and $c_j$ are pairwise distinct.
\end{claim}

If $a_i\equiv c_j\pmod p$, then $a_i+b_j\equiv 0\pmod p$, so
\[
    un+vn\equiv 0\pmod p
\]
for some $1\le u,v\le(p-1)/2$. Thus $p\mid n(u+v)$, hence $p\mid(u+v)$. But
\[
    2\le u+v\le p-1,
\]
a contradiction.

Therefore
\[
    \{a_1,\ldots,a_\ell,c_1,\ldots,c_m\}
    =
    \Bigl\{1,2,\ldots,\frac{p-1}{2}\Bigr\}.
\]
Taking products modulo $p$,
\begin{align*}
    \Bigl(\frac{p-1}{2}\Bigr)!
    &\equiv
    a_1\cdots a_\ell\,c_1\cdots c_m
    \pmod p\\
    &\equiv
    a_1\cdots a_\ell\,(p-b_1)\cdots(p-b_m)
    \pmod p\\
    &\equiv
    (-1)^m a_1\cdots a_\ell\,b_1\cdots b_m
    \pmod p\\
    &\equiv
    (-1)^m n\cdot 2n\cdots\frac{p-1}{2}n
    \pmod p\\
    &\equiv
    (-1)^m n^{(p-1)/2}\Bigl(\frac{p-1}{2}\Bigr)!
    \pmod p.
\end{align*}
The factorial is invertible modulo $p$, so
\[
    n^{(p-1)/2}\equiv(-1)^m\pmod p.
\]
Euler's criterion completes the proof.
\end{proof}

\begin{remark}[Exercise from the lecture]
Write $r_k=kn-p\bigl\lfloor kn/p\bigr\rfloor$. Counting how often $r_k>p/2$
yields the parity formula
\[
    m
    \equiv
    \frac{(n+1)(p^2-1)}{8}
    -
    \sum_{k=1}^{(p-1)/2}
    \Bigl\lfloor\frac{kn}{p}\Bigr\rfloor
    \pmod 2.
\]
In particular, if $n$ is odd (or if $p\equiv\pm 1\pmod 8$),
\[
    m
    \equiv
    \sum_{k=1}^{(p-1)/2}
    \Bigl\lfloor\frac{kn}{p}\Bigr\rfloor
    \pmod 2,
\]
while if $n$ is even and $p\equiv\pm 3\pmod 8$,
\[
    m
    \equiv
    1
    +
    \sum_{k=1}^{(p-1)/2}
    \Bigl\lfloor\frac{kn}{p}\Bigr\rfloor
    \pmod 2.
\]
The odd-numerator case will be used in the proof of quadratic reciprocity.
\end{remark}

% ============================================================
\section*{Summary of the Lecture}
% ============================================================

The lecture develops the following chain of ideas:
\[
\boxed{
\text{Wilson}
\longrightarrow
\text{quadratic residues}
\longrightarrow
\text{Legendre symbol}
\longrightarrow
\text{Euler's criterion}
\longrightarrow
\text{Gauss's lemma}
}.
\]

The key formulas are
\[
    x^2\equiv -1\pmod p
    \quad\text{solvable}
    \quad\Longleftrightarrow\quad
    p\equiv 1\pmod 4,
\]
\[
    \left(\frac np\right)
    \equiv
    n^{(p-1)/2}\pmod p,
\]
\[
    \left(\frac{-1}{p}\right)=(-1)^{(p-1)/2},
    \qquad
    \left(\frac 2p\right)=(-1)^{(p^2-1)/8},
\]
and
\[
    \left(\frac np\right)=(-1)^m
    \qquad\text{by Gauss's lemma}.
\]

\clearpage

\section*{Lecture 5}
\addcontentsline{toc}{section}{Lecture 5}

% ============================================================
\section{The Quadratic Reciprocity Law}
% ============================================================

\begin{theorem}[Quadratic reciprocity]
Let $p$ and $q$ be distinct odd primes. Then
\[
    \boxed{
    \left(\frac pq\right)
    \left(\frac qp\right)
    =
    (-1)^{(p-1)(q-1)/4}.
    }
\]
Equivalently,
\[
\left(\frac pq\right)
=
\begin{cases}
\displaystyle\left(\frac qp\right),
    & \text{if at least one of }p,q\text{ is }\equiv 1\pmod 4,\\[2mm]
\displaystyle -\left(\frac qp\right),
    & \text{if }p\equiv q\equiv 3\pmod 4.
\end{cases}
\]
\end{theorem}

\subsection{What is a reciprocity law?}

Following Wyman, a reciprocity law for a polynomial $f(x)\in\Z[x]$ is a
rule that determines the set of primes $p$ for which $f(x)\pmod p$ splits
completely (into linear factors).

For $f(x)=x^2-n$, the congruence $x^2\equiv n\pmod p$ is solvable if and
only if $f$ splits completely modulo $p$ (when $p\nmid n$). If $p\mid n$
then $f(x)\equiv x^2\pmod p$. For $p=2$,
\[
    x^2-1\equiv(x-1)^2\pmod 2.
\]
If $n\equiv a^2\pmod p$, then
\[
    x^2-n\equiv(x-a)(x+a)\pmod p;
\]
if not, then $x^2-n$ is irreducible modulo $p$.

Thus quadratic reciprocity computes $\bigl(\frac np\bigr)$ by exchanging
the roles of the two primes, which determines the splitting of $x^2-n$.

\begin{example}[$x^2-17$]
For which odd primes $p$ does $x^2-17$ split completely modulo $p$? We need
$\bigl(\frac{17}{p}\bigr)=1$. Since $17\equiv 1\pmod 4$,
\[
    \left(\frac{17}{p}\right)=\left(\frac{p}{17}\right).
\]
The squares modulo $17$ are
\[
    \{1^2,\ldots,8^2\}
    \equiv
    \{1,4,9,16,8,2,15,13\}\pmod{17},
\]
so the quadratic residues modulo $17$ are
\[
    \{1,2,4,8,9,13,15,16\}.
\]
Hence $x^2-17$ splits modulo an odd prime $p\neq 17$ if and only if
\[
    p\equiv 1,2,4,8,9,13,15,16\pmod{17}.
\]
\end{example}

\begin{example}[$x^2-11$]
Now $11\equiv 3\pmod 4$, so
\[
    \left(\frac{11}{p}\right)
    =
    (-1)^{(p-1)/2}
    \left(\frac{p}{11}\right).
\]
The quadratic residues modulo $11$ are $\{1,3,4,5,9\}$. Therefore
$\bigl(\frac{11}{p}\bigr)=1$ if and only if either
\[
    p\equiv 1\pmod 4
    \qquad\text{and}\qquad
    p\equiv 1,3,4,5,9\pmod{11},
\]
or
\[
    p\equiv 3\pmod 4
    \qquad\text{and}\qquad
    p\equiv 2,6,7,8,10\pmod{11}.
\]
The Chinese Remainder Theorem, working modulo $44$, yields the odd classes
\[
    p\equiv 1,5,7,9,19,25,35,37,39,43\pmod{44}.
\]
Thus $x^2-11$ splits completely modulo an odd prime $p\neq 11$ precisely
for these residue classes modulo $44$.
\end{example}

% ============================================================
\section{Proof of Quadratic Reciprocity}
% ============================================================

Let $p$ and $q$ be distinct odd primes. By Gauss's lemma and the counting
exercise of Lecture~4 (the numerator $q$ is odd),
\[
    \left(\frac qp\right)=(-1)^m,
    \qquad
    m\equiv\sum_{x=1}^{(p-1)/2}\Bigl\lfloor\frac{qx}{p}\Bigr\rfloor\pmod 2,
\]
and similarly
\[
    \left(\frac pq\right)=(-1)^n,
    \qquad
    n\equiv\sum_{y=1}^{(q-1)/2}\Bigl\lfloor\frac{py}{q}\Bigr\rfloor\pmod 2.
\]
It is enough to prove
\[
    m+n\equiv\frac{(p-1)(q-1)}{4}\pmod 2.
\]

Consider the function
\[
    f(x,y)=qx-py
\]
on the rectangle
\[
    X\times Y
    =
    \Bigl\{1,\ldots,\frac{p-1}{2}\Bigr\}
    \times
    \Bigl\{1,\ldots,\frac{q-1}{2}\Bigr\}.
\]
There are $(p-1)(q-1)/4$ points in this set.

The function $f$ never vanishes: $qx=py$ would force $p\mid x$, which is
impossible. It is injective: if $f(x_1,y_1)=f(x_2,y_2)$, then
\[
    q(x_1-x_2)=p(y_1-y_2),
\]
so $p\mid(x_1-x_2)$ and $q\mid(y_1-y_2)$. The bounds force
$(x_1,y_1)=(x_2,y_2)$. Thus the image of $f$ consists of
$(p-1)(q-1)/4$ distinct nonzero integers.

If $f(x,y)>0$, then $qx>py$, so $y<qx/p$. The admissible $y$ for fixed $x$
are
\[
    1\le y\le\Bigl\lfloor\frac{qx}{p}\Bigr\rfloor.
\]
Hence the number of points with $f>0$ is
\[
    \sum_{x=1}^{(p-1)/2}\Bigl\lfloor\frac{qx}{p}\Bigr\rfloor.
\]
Likewise, the number of points with $f<0$ is
\[
    \sum_{y=1}^{(q-1)/2}\Bigl\lfloor\frac{py}{q}\Bigr\rfloor.
\]
Adding these two counts,
\[
    \frac{(p-1)(q-1)}{4}
    =
    \sum_{x=1}^{(p-1)/2}\Bigl\lfloor\frac{qx}{p}\Bigr\rfloor
    +
    \sum_{y=1}^{(q-1)/2}\Bigl\lfloor\frac{py}{q}\Bigr\rfloor.
\]
The right-hand side is $m+n$ modulo $2$, which proves quadratic reciprocity.

\begin{example}
Is $219$ a quadratic residue modulo $383$? Note that $383$ is prime and
$383\equiv 3\pmod 4$, while $219=3\cdot 73$. Then
\[
    \left(\frac{219}{383}\right)
    =
    \left(\frac{3}{383}\right)
    \left(\frac{73}{383}\right).
\]
Since $3\equiv 383\equiv 3\pmod 4$,
\[
    \left(\frac{3}{383}\right)
    =
    -\left(\frac{383}{3}\right)
    =
    -\left(\frac{2}{3}\right)
    =
    -(-1)
    =1,
\]
because $383\equiv 2\pmod 3$ and $3\equiv 3\pmod 8$.

Next $73\equiv 1\pmod 4$, so
\[
    \left(\frac{73}{383}\right)
    =
    \left(\frac{383}{73}\right)
    =
    \left(\frac{18}{73}\right)
    =
    \left(\frac{2}{73}\right)
    \left(\frac{3}{73}\right)^2
    =
    \left(\frac{2}{73}\right).
\]
Since $73\equiv 1\pmod 8$, we have $\bigl(\frac{2}{73}\bigr)=1$. Therefore
\[
    \left(\frac{219}{383}\right)=1.
\]
\end{example}

% ============================================================
\section{The Jacobi Symbol}
% ============================================================

The Jacobi symbol extends the Legendre symbol to odd composite denominators.

\begin{definition}[Jacobi symbol]
Let $m>1$ be an odd integer, with prime factorization
\[
    m=\prod_{i=1}^r p_i^{\alpha_i}.
\]
For an integer $n$, the Jacobi symbol is
\[
    \boxed{
    \left(\frac nm\right)
    =
    \prod_{i=1}^r
    \left(\frac n{p_i}\right)^{\alpha_i}.
    }
\]
It is a completely multiplicative function of $n$, taking values in
$\{0,\pm 1\}$, and the factors on the right are Legendre symbols.
\end{definition}

\begin{remark}
Let $m>1$ be odd and suppose $x^2\equiv n\pmod m$ has a solution. Then
$\bigl(\frac nm\bigr)=1$ whenever $\gcdop(n,m)=1$. Indeed, if
$a^2\equiv n\pmod m$, then for every prime $p\mid m$ we have
$a^2\equiv n\pmod p$, so each Legendre factor is $1$.

This implication is one-directional: the converse fails for composite $m$.
For instance
\[
    \left(\frac{2}{9}\right)
    =
    \left(\frac{2}{3}\right)^2
    =1,
\]
but the squares modulo $9$ are $0,1,4,7$, so $x^2\equiv 2\pmod 9$ has no
solution.
\end{remark}

\begin{theorem}[Statements used without proof]
Let $m,n>1$ be odd integers.

\begin{enumerate}[label=\arabic*.]
    \item
    \[
        \left(\frac{-1}{m}\right)=(-1)^{(m-1)/2},
        \qquad
        \left(\frac{2}{m}\right)=(-1)^{(m^2-1)/8}.
    \]

    \item If $\gcdop(m,n)=1$, then
    \[
        \left(\frac nm\right)
        \left(\frac mn\right)
        =
        (-1)^{(m-1)(n-1)/4}.
    \]
\end{enumerate}
\end{theorem}

% ============================================================
\section{Polynomial Congruences Modulo Prime Powers}
% ============================================================

By the Chinese Remainder Theorem, solving $f(x)\equiv 0\pmod m$ reduces to
solving
\[
    f(x)\equiv 0\pmod{p^\alpha}
\]
for every prime power $p^\alpha$ exactly dividing $m$. A solution modulo
$p^\alpha$ automatically gives a solution modulo $p^\beta$ for every
$1\le\beta\le\alpha$: if $f(a)\equiv 0\pmod{p^\alpha}$ and $\beta<\alpha$,
write $a=kp^\beta+r$ with $0\le r<p^\beta$; then $a\equiv r\pmod{p^\beta}$,
so $f(r)\equiv f(a)\equiv 0\pmod{p^\beta}$.

The converse lifting problem is: given $0\le r<p^{\alpha-1}$ with
$f(r)\equiv 0\pmod{p^{\alpha-1}}$, can one find $0\le a<p^\alpha$ with
\[
    f(a)\equiv 0\pmod{p^\alpha}
    \qquad\text{and}\qquad
    a\equiv r\pmod{p^{\alpha-1}}?
\]

\begin{theorem}[Hensel's lemma]
Let $p$ be a prime, let $\alpha\ge 2$, and let $f\in\Z[x]$. Suppose
$0\le r<p^{\alpha-1}$ satisfies
\[
    f(r)\equiv 0\pmod{p^{\alpha-1}}.
\]

\begin{enumerate}[label=\arabic*.]
    \item If $f'(r)\not\equiv 0\pmod p$, then there exists a unique $a$ with
    $0\le a<p^\alpha$ such that
    \[
        f(a)\equiv 0\pmod{p^\alpha}
        \qquad\text{and}\qquad
        a\equiv r\pmod{p^{\alpha-1}}.
    \]

    \item If $f'(r)\equiv 0\pmod p$, then:
    \begin{enumerate}[label=(\roman*)]
        \item if $f(r)\equiv 0\pmod{p^\alpha}$, there are exactly $p$ distinct
        solutions $a$ modulo $p^\alpha$ with $a\equiv r\pmod{p^{\alpha-1}}$;

        \item if $f(r)\not\equiv 0\pmod{p^\alpha}$, there is no such solution.
    \end{enumerate}
\end{enumerate}
\end{theorem}

\begin{proof}[Proof (not given in lecture)]
Every candidate lift has the form $a=r+tp^{\alpha-1}$ with $0\le t<p$.
Taylor expansion gives
\[
    f(r+tp^{\alpha-1})
    =
    f(r)+f'(r)\,tp^{\alpha-1}
    +\frac{f''(r)}{2}(tp^{\alpha-1})^2
    +\cdots.
\]
For $\alpha\ge 2$ one has $2(\alpha-1)\ge\alpha$, so every term of degree
$\ge 2$ is $0$ modulo $p^\alpha$. Thus
\[
    f(r+tp^{\alpha-1})
    \equiv
    f(r)+f'(r)\,tp^{\alpha-1}
    \pmod{p^\alpha}.
\]
Write $f(r)=kp^{\alpha-1}$. The congruence $f(a)\equiv 0\pmod{p^\alpha}$
becomes
\[
    k+f'(r)\,t\equiv 0\pmod p.
\]
If $f'(r)\not\equiv 0\pmod p$, there is a unique solution $t$ modulo $p$.
If $f'(r)\equiv 0\pmod p$, then either $k\equiv 0\pmod p$, in which case
every $t$ works, or $k\not\equiv 0\pmod p$, in which case there is no
solution.
\end{proof}

\begin{example}
The congruence $x^2\equiv 2\pmod 4$ has no solution, since squares modulo
$4$ are only $0$ and $1$. Thus a solution modulo $p$ need not lift when
the derivative condition fails.
\end{example}

% ============================================================
\section*{Summary of the Lecture}
% ============================================================

The lecture develops the following chain of ideas:
\[
\boxed{
\text{Gauss's lemma}
\longrightarrow
\text{quadratic reciprocity}
\longrightarrow
\text{Jacobi symbol}
\longrightarrow
\text{Hensel lifting}
}.
\]

The key formulas are
\[
    \left(\frac pq\right)
    \left(\frac qp\right)
    =
    (-1)^{(p-1)(q-1)/4},
\]
the Jacobi symbol extends the Legendre symbol (though
$\bigl(\tfrac nm\bigr)=1$ need not imply solvability for composite $m$),
and Hensel gives a unique lift when $f'(r)\not\equiv 0\pmod p$.

\clearpage

\section*{Lecture 6}
\addcontentsline{toc}{section}{Lecture 6}

% ============================================================
\section{Congruences Modulo Prime Powers}
% ============================================================

A solution of $f(x)\equiv 0\pmod{p^\alpha}$ with $\alpha\ge 2$ is in
particular a solution modulo $p^{\alpha-1}$. Hensel's lemma of the previous
lecture addresses the converse: when a solution modulo $p^{\alpha-1}$ lifts
to a solution modulo $p^\alpha$ in the same residue class.

% ============================================================
\section{Units and Primitive Roots}
% ============================================================

Recall that a complete residue system modulo $m$ represents $\Z/m\Z$, while
a reduced residue system represents the unit group
\[
    \units m.
\]
For example,
\[
    \units{2}=\{\overline 1\},
    \qquad
    \units{3}=\{\overline 1,\overline 2\},
    \qquad
    \units{4}=\{\overline 1,\overline 3\}.
\]
Groups of prime order are cyclic. Direct inspection gives:
\begin{align*}
    m=5&\colon
    \units 5=\{\overline 1,\overline 2,\overline 3,\overline 4\}
    \text{ is cyclic of order }4,\\
    m=6&\colon
    \units 6=\{\overline 1,\overline 5\}
    \text{ is cyclic},\\
    m=7&\colon
    \units 7
    \text{ is cyclic, with generator }\overline 5,\\
    m=8&\colon
    \units 8=\{\overline 1,\overline 3,\overline 5,\overline 7\}
    \text{ is abelian but not cyclic}.
\end{align*}

By Euler's theorem, every $a\in\units m$ satisfies
\[
    a^{\phiE(m)}\equiv 1\pmod m.
\]
The natural question is whether there exists $a$ whose order is exactly
$\phiE(m)$, i.e.\ $a^{\phiE(m)}\equiv 1\pmod m$ but $a^k\not\equiv 1\pmod m$
for every $0<k<\phiE(m)$. In that case $\units m$ is cyclic.

The complex analogue is a primitive $a$-th root of unity: a solution of
$X^a=1$ which is not a solution of $X^b=1$ for any $0<b<a$. There are
$\phiE(a)$ primitive $a$-th roots of unity.

The multiplicative order $\ord_m(a)$ of a unit $a$ divides $\phiE(m)$.

\begin{definition}[Primitive root]
Let $m>1$ and let $a\in\Z$ with $\gcdop(a,m)=1$. We say that $a$ is a
\emph{primitive root modulo $m$} if
\[
    \ord_m(a)=\phiE(m).
\]
Equivalently,
\[
    \{a,a^2,\ldots,a^{\phiE(m)}\}
    =
    \units m
\]
as sets of residue classes. Thus $a$ is a primitive root modulo $m$ if and
only if $\units m$ is cyclic and $a$ generates it.
\end{definition}

\begin{remark}[Exercise]
Let $m>1$, $\gcdop(a,m)=1$, and $\ord_m(a)=k$. Then:

\begin{enumerate}[label=\arabic*.]
    \item $a^g\equiv a^h\pmod m$ if and only if $g\equiv h\pmod k$;

    \item $a^g\equiv 1\pmod m$ if and only if $k\mid g$;

    \item $1,a,a^2,\ldots,a^{k-1}$ are pairwise incongruent modulo $m$.
\end{enumerate}
\end{remark}

% ============================================================
\section{The Unit Group Modulo a Prime is Cyclic}
% ============================================================

\begin{lemma}
Let $m>1$, $\gcdop(a,m)=1$, and $\ord_m(a)=k$. Then
\[
    \ord_m(a^g)
    =
    \frac{k}{\gcdop(k,g)}.
\]
\end{lemma}

\begin{proof}
Let $t=\ord_m(a^g)$. Then $a^{gt}\equiv 1\pmod m$, so $k\mid gt$, or
equivalently $k/\gcdop(k,g)$ divides $t$.

Conversely,
\[
    (a^g)^{k/\gcdop(k,g)}=a^{k\cdot(g/\gcdop(k,g))}\equiv 1,
\]
so $t$ divides $k/\gcdop(k,g)$.
\end{proof}

\begin{lemma}
For every integer $n\ge 1$,
\[
    \sum_{d\mid n}\phiE(d)=n.
\]
\end{lemma}

\begin{proof}
Partition $\{1,2,\ldots,n\}$ according to $\gcdop(k,n)$. For each $d\mid n$,
set
\[
    A_d=\{k:1\le k\le n,\ \gcdop(k,n)=d\}.
\]
The map $k\mapsto k/d$ is a bijection
\[
    A_d
    \longrightarrow
    \{j:1\le j\le n/d,\ \gcdop(j,n/d)=1\},
\]
so $|A_d|=\phiE(n/d)$. Therefore
\[
    n=\sum_{d\mid n}|A_d|=\sum_{d\mid n}\phiE(n/d)=\sum_{c\mid n}\phiE(c).
\]
\end{proof}

\begin{theorem}
Let $p$ be a prime and let $d\mid(p-1)$. Among $1,2,\ldots,p-1$ there are
exactly $\phiE(d)$ integers of multiplicative order $d$ modulo $p$.
\end{theorem}

\begin{corollary}
Taking $d=p-1$, there are $\phiE(p-1)$ elements of order $p-1$. In
particular $\units p$ is cyclic, and there exists a primitive root modulo
every prime $p$.
\end{corollary}

\begin{proof}
For each $d\mid(p-1)$ define
\[
    B_d=\{a:1\le a\le p-1,\ \ord_p(a)=d\}.
\]
These sets partition the nonzero residues, so
\[
    p-1=\sum_{d\mid(p-1)}|B_d|.
\]

\begin{claim}
If $B_d\neq\varnothing$, then $|B_d|=\phiE(d)$.
\end{claim}

Suppose $a\in B_d$, so $a^d\equiv 1\pmod p$ and $\ord_p(a)=d$. The $d$
powers $a,a^2,\ldots,a^d$ are distinct and all satisfy $X^d\equiv 1\pmod p$.
The polynomial $X^d-1$ has at most $d$ roots in the field $\Z/p\Z$, so these
powers are \emph{all} the solutions of $X^d\equiv 1$. Any element of order
$d$ is among them. By the first lemma,
\[
    \ord_p(a^k)=d
    \quad\Longleftrightarrow\quad
    \gcdop(k,d)=1.
\]
There are exactly $\phiE(d)$ such exponents, so $|B_d|=\phiE(d)$.

Therefore
\[
    p-1
    =
    \sum_{\substack{d\mid(p-1)\\ B_d\neq\varnothing}}\phiE(d)
    \le
    \sum_{d\mid(p-1)}\phiE(d)
    =p-1.
\]
Equality throughout forces $B_d\neq\varnothing$ for every $d\mid(p-1)$. In
particular $B_{p-1}\neq\varnothing$, so a primitive root modulo $p$ exists,
and $|B_d|=\phiE(d)$ for every such $d$.
\end{proof}

Thus for a prime $p$ there exists $a\in\{1,\ldots,p-1\}$ with
\[
    \{a,a^2,\ldots,a^{p-1}\}=\{1,2,\ldots,p-1\}
\]
as residue classes modulo $p$.

\begin{remark}
There is no integer $a$ which is a primitive root modulo every prime
(consider a prime factor of $a$). There is also no integer $a$ which is a
primitive root modulo all but finitely many primes.
\end{remark}

% ============================================================
\section{Primitive Roots Cannot Be Squares}
% ============================================================

\begin{proposition}
Let $p$ be an odd prime. A primitive root modulo $p$ cannot be a quadratic
residue modulo $p$.
\end{proposition}

\begin{proof}
If $a\equiv b^2\pmod p$ with $p\nmid a$, Euler's criterion gives
\[
    a^{(p-1)/2}\equiv 1\pmod p,
\]
so $\ord_p(a)$ divides $(p-1)/2$ and cannot equal $p-1$.
\end{proof}

\begin{theorem}
Given an integer $a\in\Z$, there are infinitely many primes $p$ such that
$a$ is not a primitive root modulo $p$.
\end{theorem}

\begin{proof}
If $a=0$, then $\gcdop(0,p)=p\neq1$ for every prime $p$, so $0$ is never a
primitive root. Thus every prime works.

Now assume $a\neq 0$. It is enough to produce infinitely many primes modulo
which $a$ is a nonzero square (then the previous proposition applies for
odd primes, while $p=2$ is at most one exception). Write
\[
    a=\pm\,2^e b^2 c
\]
with $c$ odd and square-free, $c=p_1\cdots p_k$ a product of distinct odd
primes (the finitely many primes dividing $a$ may be ignored).

If $p\equiv 1\pmod 4$ and $p\equiv 1\pmod{p_i}$ for each $i$, quadratic
reciprocity gives
\[
    \left(\frac{p_i}{p}\right)
    =
    \left(\frac{p}{p_i}\right)
    =
    \left(\frac{1}{p_i}\right)
    =1,
\]
and likewise $\bigl(\frac{-1}{p}\bigr)=1$. The symbol $\bigl(\frac{2}{p}\bigr)$
can be controlled by a congruence modulo $8$ if needed.

By the Chinese Remainder Theorem there is a residue class modulo
$4p_1\cdots p_k$ (or modulo $8$ times that product) consisting of integers
$n\equiv 1$ at each of these moduli. Dirichlet's theorem on primes in
arithmetic progressions supplies infinitely many primes in that class, and
for all such primes $\bigl(\frac ap\bigr)=1$.
\end{proof}

\subsection{Artin's primitive root conjecture}

A modified question: does there exist an integer $a$ which is a primitive
root for infinitely many primes $p$?

A stronger question: if $a$ is not a square and $a\neq -1$, do there exist
infinitely many primes $p$ for which $a$ is a primitive root modulo $p$?

Artin's primitive root conjecture asserts that the answer to the stronger
question is yes. The conjecture is still open; one does not even know that
the modified question has a positive answer for some explicit $a$.

Known partial results recorded in the lecture:

\begin{itemize}
    \item Heath--Brown: at least one of $2$, $3$, $5$ is a primitive root for
    infinitely many primes.

    \item Hooley: Artin's conjecture is true if the Generalized Riemann
    Hypothesis is assumed.
\end{itemize}

% ============================================================
\section{Which Moduli Admit Primitive Roots?}
% ============================================================

\begin{theorem}
Let $m\ge 2$. There exists a primitive root modulo $m$ if and only if
\[
    m\in\{2,4,p^n,2p^n\},
\]
where $p$ is an odd prime and $n\ge 1$. Equivalently, $\units m$ is cyclic
precisely for these $m$.
\end{theorem}

The cases $m=2,4$ and the prime case $m=p$ have been checked. The remaining
prime-power and twice prime-power cases, and the failure for all other $m$
(already visible for $m=8$), are the content of the standard structure
theorem for $\units m$.

% ============================================================
\section*{Summary of the Lecture}
% ============================================================

The lecture develops the following chain of ideas:
\[
\boxed{
\text{unit groups}
\longrightarrow
\text{primitive roots}
\longrightarrow
(\Z/p\Z)^\times\text{ cyclic}
\longrightarrow
\text{Artin conjecture}
}.
\]

The main results are that $\units p$ is cyclic for every prime $p$, a
primitive root modulo an odd prime cannot be a square, and $\units m$ is
cyclic if and only if $m\in\{2,4,p^n,2p^n\}$.

\clearpage

\section*{Lecture 7}
\addcontentsline{toc}{section}{Lecture 7}

% ============================================================
\section{Primitive Roots Modulo $m$}
% ============================================================

Write $S$ for the set of moduli which admit a primitive root:
\[
    S=\{2,4,p^k,2p^k:p\text{ an odd prime},\ k\ge 1\}.
\]

\begin{theorem}
The group $\units m$ is cyclic if and only if $m\in S$. In that case there
are exactly
\[
    \boxed{\phiE(\phiE(m))}
\]
primitive roots modulo $m$. Equivalently, these are the residue classes $x$
satisfying $\ord_m(x)=\phiE(m)$, and
\[
    \units m=\{a,a^2,\ldots,a^{\phiE(m)}\}
\]
whenever $a$ is one such primitive root.
\end{theorem}

\subsection{A Useful Criterion}

Suppose $n>1$. If there exists an integer $a\ge1$ such that
\[
    a^{n-1}\equiv1\pmod n
\]
and, for every prime divisor $q$ of $n-1$,
\[
    a^{(n-1)/q}\not\equiv1\pmod n,
\]
then $n$ is prime.

Indeed, the order of $a$ modulo $n$ must divide $n-1$. The displayed
conditions force
\[
    \ord_n(a)=n-1.
\]
Since
\[
    \ord_n(a)\mid\phiE(n),
\]
we obtain
\[
    n-1\mid\phiE(n).
\]
But
\[
    \phiE(n)\le n-1
\]
for $n>1$, and hence
\[
    \phiE(n)=n-1.
\]
Therefore $n$ is prime.

% ============================================================
\section{Index / Discrete Logarithm}
% ============================================================

Let $m\in\N$ be such that $\units m$ is cyclic, and let $a$ be a
primitive root modulo $m$. Then
\[
    \units m
    =
    \{a,a^2,\ldots,a^{\phiE(m)}\}.
\]

\begin{definition}[Index / discrete logarithm]
For any
\[
    b\in\units m,
\]
there exists a unique integer
\[
    1\le n\le\phiE(m)
\]
such that
\[
    b\equiv a^n\pmod m.
\]
We define
\[
    \boxed{\ind_a b=n}
\]
and call $n$ the \emph{index} (or discrete logarithm) of $b$ to the base
$a$ modulo $m$.
\end{definition}

Thus
\[
    b\equiv a^{\ind_a b}\pmod m.
\]

For positive real numbers $x,y,z$ with $x\neq 1$ and $y\neq 1$,
\[
    x=y^z
    \quad\Longrightarrow\quad
    \log_y x=z,
\]
and hence
\[
    \log_x z
    =
    \frac{\log_y z}{\log_y x}.
\]
The index satisfies analogous change-of-base and logarithm-like identities.

\subsection{Basic Properties of the Index}

Let $m\in S$, where $S$ denotes the moduli for which a primitive root
exists, and let $a$ be a primitive root modulo $m$.

Let $b_1,b_2$ be integers satisfying
\[
    \gcdop(b_1,m)=\gcdop(b_2,m)=1.
\]
Then:

\begin{enumerate}[label=\Roman*.]
    \item
    \[
        \ind_a(b_1b_2)
        \equiv
        \ind_a b_1+\ind_a b_2
        \pmod{\phiE(m)}.
    \]

    \item For any integer $n$,
    \[
        \ind_a(b_1^n)
        \equiv
        n\,\ind_a b_1
        \pmod{\phiE(m)}.
    \]

    \item
    \[
        \ind_a(1)
        \equiv
        \phiE(m)
        \equiv0
        \pmod{\phiE(m)}.
    \]

    \item If $m>2$, then
    \[
        \ind_a(-1)
        \equiv
        \frac{\phiE(m)}2
        \pmod{\phiE(m)}.
    \]

    \item Let $a,a_1$ be two primitive roots modulo $m$ and $b\in\units m$.
    Then
    \[
        \ind_a b
        \equiv
        (\ind_{a_1} b)\,(\ind_a a_1)
        \pmod{\phiE(m)},
    \]
    the analogue of $\log_x y=(\log_z y)(\log_x z)$.
\end{enumerate}

\begin{proof}[Proof (not given in lecture)]
Write
\[
    b_1\equiv a^{r}\pmod m,
    \qquad
    b_2\equiv a^{s}\pmod m.
\]
Then
\[
    b_1b_2\equiv a^{r+s}\pmod m,
\]
which gives the first identity modulo $\phiE(m)$.

Similarly,
\[
    b_1^n\equiv a^{nr}\pmod m,
\]
giving the second identity.

Since
\[
    a^{\phiE(m)}\equiv1\pmod m,
\]
the index of $1$ is $0$ modulo $\phiE(m)$, which is represented by
$\phiE(m)$ in the lecture's convention.

Finally,
\[
    (-1)^2\equiv1\pmod m,
\]
and the order of $-1$ is $2$ (for $m>2$). Therefore its index is the
unique exponent giving order $2$, namely
\[
    \frac{\phiE(m)}2
\]
modulo $\phiE(m)$.

For the change-of-base formula, write $b\equiv a_1^{r}\pmod m$ and
$a_1\equiv a^{s}\pmod m$. Then $b\equiv a^{rs}\pmod m$, so
$\ind_a b\equiv rs\equiv(\ind_{a_1}b)(\ind_a a_1)$ modulo $\phiE(m)$.
\end{proof}

\subsection{Squares via the Index}

Let $m\in S$ with $m>2$, and let $b\in\units m$. Then $b$ is a square modulo
$m$ if and only if $\ind_a b$ is even, for some primitive root $a$ modulo
$m$. Equivalently, the index is even for every primitive root modulo $m$.

\subsection{Solving Exponential Congruences}

Consider
\[
    x^k\equiv b\pmod m,
    \qquad
    m\in S,
    \qquad
    k\ge 1,
    \qquad
    \gcdop(b,m)=1,
\]
where $x$ is sought in $\units m$.
Let $a$ be a primitive root modulo $m$. Write
\[
    b\equiv a^{\ind_a b}\pmod m
\]
and
\[
    x\equiv a^{\ind_a x}\pmod m.
\]
Then
\[
    x^k\equiv b\pmod m
\]
is equivalent to
\[
    k\,\ind_a x
    \equiv
    \ind_a b
    \pmod{\phiE(m)}.
\]

Therefore the exponential congruence has a solution if and only if
\[
    \boxed{
    \gcdop\!\left(k,\phiE(m)\right)
    \mid
    \ind_a b
    }.
\]

In that case, it has exactly
\[
    \gcdop\!\left(k,\phiE(m)\right)
\]
solutions.

\begin{theorem}
For $m\in S$, $k\ge 1$, and $\gcdop(b,m)=1$, the congruence
\[
    x^k\equiv b\pmod m
\]
in $\units m$ has a solution if and only if
\[
    \boxed{
    b^{\phiE(m)/d}\equiv1\pmod m
    },
    \qquad
    d=\gcdop(k,\phiE(m)).
    \]
\end{theorem}

\begin{proof}
Since $a$ is primitive,
\[
    b\equiv a^{\ind_a b}\pmod m.
\]
The linear congruence
\[
    k\,\ind_a x
    \equiv
    \ind_a b
    \pmod{\phiE(m)}
\]
is solvable if and only if
\[
    d\mid\ind_a b,
    \qquad
    d=\gcdop(k,\phiE(m)).
\]
This is equivalent to
\[
    \phiE(m)\mid
    \frac{\phiE(m)}d\,\ind_a b.
\]
Hence
\[
    a^{(\phiE(m)/d)\ind_a b}
    \equiv1\pmod m.
\]
But
\[
    a^{\ind_a b}\equiv b\pmod m,
\]
so this is exactly
\[
    b^{\phiE(m)/d}\equiv1\pmod m.
\]
\end{proof}

\subsection{Example: $x^3\equiv4\pmod{13}$}

Consider
\[
    x^3\equiv4\pmod{13}.
\]
Take $2$ as a primitive root modulo $13$. The powers of $2$ modulo $13$
are
\[
\begin{array}{c|cccccccccccc}
k&1&2&3&4&5&6&7&8&9&10&11&12\\
\hline
2^k&2&4&8&3&6&12&11&9&5&10&7&1
\end{array}
\]
so
\[
    \ind_2(4)=2.
\]
The congruence becomes
\[
    3\ind_2 x\equiv2\pmod{12}.
\]
Since
\[
    \gcdop(3,12)=3
\]
and
\[
    3\nmid2,
\]
there is no solution.

Thus
\[
    \boxed{x^3\equiv4\pmod{13}\text{ has no solution}.}
\]

\subsection{Example: $x^3\equiv5\pmod{13}$}

Again use $2$ as a primitive root modulo $13$. Since
\[
    2^9\equiv5\pmod{13},
\]
we have
\[
    \ind_2(5)=9.
\]
Thus
\[
    3\ind_2 x\equiv9\pmod{12}.
\]
Dividing by $3$ gives
\[
    \ind_2 x\equiv3\pmod4.
\]
Therefore
\[
    \ind_2 x\in\{3,7,11\}.
\]
The corresponding solutions are
\[
    x\in\{2^3,2^7,2^{11}\}\pmod{13},
\]
i.e.
\[
    \boxed{x\equiv8,11,7\pmod{13}.}
\]

\subsection{Example: $4x^9\equiv 7\pmod{13}$}

Again take the primitive root $2$ modulo $13$. From the table,
\[
    \ind_2 4=2,
    \qquad
    \ind_2 7=11.
\]
The congruence becomes
\[
    \ind_2 4+9\ind_2 x\equiv\ind_2 7\pmod{12},
\]
that is
\[
    2+9\ind_2 x\equiv 11\pmod{12},
    \qquad
    9\ind_2 x\equiv 9\pmod{12}.
\]
Dividing by $3$ yields
\[
    3\ind_2 x\equiv 3\pmod 4,
    \qquad
    \ind_2 x\equiv 1\pmod 4.
\]
Thus $\ind_2 x\in\{1,5,9\}$, and
\[
    \boxed{x\equiv 2,\ 6,\ 5\pmod{13}.}
\]

% ============================================================
\section{Exponential Congruences: General Form}
% ============================================================

Consider
\[
    b^x\equiv c\pmod m,
    \qquad
    m\in S,
    \qquad
    \gcdop(b,m)=\gcdop(c,m)=1.
\]
Taking indices with respect to a primitive root $a$ gives
\[
    \ind_a(b^x)
    \equiv
    \ind_a(c)
    \pmod{\phiE(m)}.
\]
Using the index laws,
\[
    x\,\ind_a b
    \equiv
    \ind_a c
    \pmod{\phiE(m)}.
\]

Thus the problem reduces to a linear congruence:
\[
    \boxed{
    x\,\ind_a b
    \equiv
    \ind_a c
    \pmod{\phiE(m)}
    }.
\]

It has a solution precisely when
\[
    \boxed{
    \gcdop(\ind_a b,\phiE(m))
    \mid
    \ind_a c
    }.
\]

\begin{example}
As an exercise, solve
\[
    10^x\equiv 9\pmod{13}.
\]
\end{example}

% ============================================================
\section{Finite Fields}
% ============================================================

\subsection{The Field $\Z/2\Z$}

The lecture begins the finite-field discussion with
\[
    \Z/2\Z.
\]

Let $F$ be a finite field with $q$ elements. Then
\[
    F^\times=F\setminus\{0\}
\]
is a multiplicative group with
\[
    |F^\times|=q-1.
\]
Consequently every element of $F^\times$ satisfies
\[
    x^{q-1}=1.
\]
Equivalently, every element of $F$ satisfies
\[
    \boxed{x^q-x=0.}
\]

\subsection{Polynomial Root Bound}

\begin{lemma}
Let $F$ be a field and let
\[
    f(x)\in F[x]
\]
be a non-zero polynomial of degree $n\ge1$. Then $f$ has at most $n$
distinct roots in $F$.
\end{lemma}

\begin{proof}
Suppose $f$ has distinct roots
\[
    \alpha_1,\ldots,\alpha_r.
\]
Then each root gives a factor:
\[
    (x-\alpha_i)\mid f(x).
\]
Since the roots are distinct,
\[
    \prod_{i=1}^r(x-\alpha_i)\mid f(x).
\]
Therefore
\[
    r\le\deg f=n.
\]
\end{proof}

\subsection{Factorization of $x^q-x$}

Since every element of $F$ satisfies
\[
    x^q-x=0,
\]
the polynomial $x^q-x$ has all $q$ elements of $F$ as roots. Since its
degree is $q$, it follows that
\[
    \boxed{
    x^q-x
    =
    \prod_{\alpha\in F}(x-\alpha)
    }.
\]

\begin{corollary}
Let $F$ be a finite field with $q$ elements, and let $L$ be any extension
field of $F$. An element $\alpha\in L$ lies in $F$ if and only if
$\alpha^q=\alpha$.
\end{corollary}

\begin{corollary}
Let $F$ be a finite field with $q$ elements and let $f(x)\in F[x]$ have
degree $d$ and divide $x^q-x$. Then $f$ has $d$ distinct roots in $F$.
\end{corollary}

\begin{proof}
Write $x^q-x=f(x)g(x)$. The left-hand side is a product of $q$ distinct
linear factors, so $f$ is square-free and splits completely in $F$.
\end{proof}

More generally, if $F$ is a finite field with $q$ elements and
$K\subseteq F$ is a subfield with $|K|=q'$, then every $\alpha\in K$
satisfies $\alpha^{q'}=\alpha$.

% ============================================================
\section{Cyclicity of the Multiplicative Group of a Finite Field}
% ============================================================

Finite fields share the cyclicity property observed earlier for the
groups
\[
    \units m
\]
for the moduli $m$ admitting primitive roots.

\begin{theorem}
Let $F$ be a finite field with $q$ elements. Then
\[
    \boxed{F^\times\text{ is cyclic}.}
\]
In particular, there exists $\alpha\in F^\times$ such that
\[
    \ord(\alpha)=q-1.
\]
\end{theorem}

\subsection{Proof}

Every element $\alpha\in F^\times$ satisfies
\[
    \alpha^{q-1}=1.
\]
Hence the order of every element divides $q-1$.

For every divisor
\[
    d\mid q-1,
\]
define
\[
    B_d
    =
    \{\,\alpha\in F^\times:\ord(\alpha)=d\,\}.
\]

The sets $B_d$ partition $F^\times$. Thus
\[
    q-1
    =
    \sum_{d\mid q-1}|B_d|.
\]

\subsection{Claim: $|B_d|\le\phiE(d)$}

Suppose
\[
    B_d\neq\varnothing
\]
and choose $\alpha\in B_d$. Then
\[
    \ord(\alpha)=d.
\]
The elements
\[
    \alpha,\alpha^2,\ldots,\alpha^d
\]
are all distinct, and all satisfy $X^d=1$. The polynomial $X^d-1$ has at
most $d$ roots in the field $F$, so these powers are all the solutions of
$X^d=1$. Any element of order $d$ is therefore among them.

For
\[
    1\le k\le d,
\]
we have
\[
    \ord(\alpha^k)
    =
    \frac{d}{\gcdop(k,d)}.
\]
Therefore
\[
    \ord(\alpha^k)=d
    \quad\Longleftrightarrow\quad
    \gcdop(k,d)=1.
\]
There are exactly $\phiE(d)$ such values of $k$. Hence
\[
    |B_d|=\phiE(d)
\]
whenever $B_d\neq\varnothing$.

Consequently,
\[
    q-1
    =
    \sum_{\substack{d\mid q-1\\B_d\neq\varnothing}}
    \phiE(d).
\]

\subsection{Existence of an Element of Order $q-1$}

Suppose, for contradiction, that $F^\times$ has no element of order
$q-1$. Then
\[
    B_{q-1}=\varnothing.
\]
Hence
\[
    q-1
    =
    \sum_{\substack{d\mid q-1\\d<q-1\\B_d\neq\varnothing}}
    \phiE(d)
    \le
    \sum_{\substack{d\mid q-1\\d<q-1}}\phiE(d).
\]
But
\[
    \sum_{d\mid n}\phiE(d)=n,
\]
so
\[
    \sum_{\substack{d\mid q-1\\d<q-1}}\phiE(d)
    =
    (q-1)-\phiE(q-1)
    <q-1,
\]
a contradiction.

Therefore
\[
    B_{q-1}\neq\varnothing.
\]
Thus there exists $\alpha\in F^\times$ with
\[
    \ord(\alpha)=q-1,
\]
and hence
\[
    F^\times=\langle\alpha\rangle.
\]

This proves that
\[
    \boxed{F^\times\text{ is cyclic}.}
\]

% ============================================================
\section{Finite Fields Have Prime-Power Size}
% ============================================================

Let $F$ be a finite field with $q$ elements.

Consider the prime subfield
\[
    K=\{0,1,2,\ldots,p-1\},
\]
where $p$ is the characteristic of $F$. Thus
\[
    K\cong\Z/p\Z
\]
and $K$ has $p$ elements.

Hence $K$ is a subfield of $F$:
\[
    K\subseteq F.
\]

Now regard $F$ as a vector space over $K$. Let
\[
    \dim_KF=n.
\]
Choose a basis
\[
    e_1,\ldots,e_n.
\]
Every element of $F$ has a unique representation
\[
    a_1e_1+\cdots+a_ne_n,
    \qquad
    a_i\in K.
\]
Since there are $p$ choices for each coefficient,
\[
    |F|=p^n.
\]

Therefore every finite field has prime-power cardinality:
\[
    \boxed{|F|=p^n.}
\]

In particular, there are no finite fields with $6$ elements, since $6$ is
not a prime power.

\begin{remark}
The lecture notes that a finite field of size $q$ must have
\[
    q=p^n
\]
for a prime $p$ and integer $n\ge1$.
\end{remark}

% ============================================================
\section*{Summary}
% ============================================================

The lecture develops the following chain:
\[
\boxed{
\text{Primitive roots}
\longrightarrow
\text{cyclic }\units m
\longrightarrow
\text{index / discrete logarithm}
\longrightarrow
\text{solving exponential congruences}
}
\]
and then moves to finite fields:
\[
\boxed{
\text{finite field}
\longrightarrow
F^\times
\longrightarrow
\text{orders of elements}
\longrightarrow
F^\times\text{ cyclic}
}.
\]

The main formulas and results are
\[
    \ord_m(a)=\phiE(m)
    \quad\Longleftrightarrow\quad
    a\text{ is a primitive root modulo }m,
\]
\[
    \units m\text{ is cyclic}
    \quad\Longleftrightarrow\quad
    m\in\{2,4,p^k,2p^k\},
\]
\[
    \#\{\text{primitive roots mod }m\}
    =
    \phiE(\phiE(m))
    \qquad
    \text{when }\units m\text{ is cyclic},
\]
\[
    \ind_a(b_1b_2)
    \equiv
    \ind_a b_1+\ind_a b_2
    \pmod{\phiE(m)},
\]
\[
    x^k\equiv b\pmod m
    \quad\Longleftrightarrow\quad
    k\ind_a x\equiv\ind_a b
    \pmod{\phiE(m)},
\]
\[
    x^k\equiv b\pmod m
    \text{ is solvable}
    \quad\Longleftrightarrow\quad
    b^{\phiE(m)/d}\equiv1\pmod m,
    \quad
    d=\gcdop(k,\phiE(m)),
\]
\[
    x^q-x
    =
    \prod_{\alpha\in F}(x-\alpha)
    \qquad(|F|=q),
\]
\[
    \boxed{F^\times\text{ is cyclic for every finite field }F},
\]
and
\[
    \boxed{|F|=p^n}
\]
for some prime $p$ and integer $n\ge1$.

\clearpage

\section*{Lecture 8}
\addcontentsline{toc}{section}{Lecture 8}

% ============================================================
\section{Finite Fields and Their Characteristic}
% ============================================================

Let $F$ be a finite field with $q$ elements.

A finite field with $q$ elements is characterized by the equation
\[
    \boxed{x^q-x=0.}
\]
Indeed, every element $\alpha\in F$ satisfies
\[
    \alpha^q=\alpha.
\]

\begin{proposition}
The additive order of $1\in F$ is a prime number $p$, called the
characteristic of $F$.
\end{proposition}

Thus
\[
    \charp(F)=p
\]
for some prime $p$.

The prime subfield of $F$ is
\[
    \boxed{\Z/p\Z}.
\]

% ============================================================
\section{Freshman's Dream / Wishful Shortcut}
% ============================================================

\begin{theorem}[Freshman's dream / wishful shortcut]
Let $F$ be a finite field of characteristic $p$. Then for
\[
    \alpha,\beta\in F
\]
and
\[
    n\in\N,
\]
we have
\[
    \boxed{
    (\alpha+\beta)^{p^n}
    =
    \alpha^{p^n}+\beta^{p^n}.
    }
\]
\end{theorem}

\begin{proof}
For $n=1$,
\[
    (\alpha+\beta)^p
    =
    \sum_{i=0}^{p}
    \binom{p}{i}\alpha^i\beta^{p-i}.
\]
For
\[
    1\le i\le p-1,
\]
we have
\[
    p\mid\binom{p}{i}.
\]
Since $\charp(F)=p$, all the intermediate binomial coefficients vanish
in $F$. Hence
\[
    (\alpha+\beta)^p
    =
    \alpha^p+\beta^p.
\]

Iterating this identity gives
\[
    (\alpha+\beta)^{p^2}
    =
    \left((\alpha+\beta)^p\right)^p
    =
    (\alpha^p+\beta^p)^p
    =
    \alpha^{p^2}+\beta^{p^2},
\]
and in general
\[
    (\alpha+\beta)^{p^n}
    =
    \alpha^{p^n}+\beta^{p^n}.
\]
\end{proof}

% ============================================================
\section{Divisibility of $x^\ell-1$ and $x^m-1$}
% ============================================================

\begin{theorem}
Let $F$ be a field, and let $\ell,m$ be positive integers. Then
\[
    \boxed{
    x^\ell-1\divides x^m-1
    \quad\Longleftrightarrow\quad
    \ell\mid m.
    }
\]
\end{theorem}

\begin{proof}
First suppose
\[
    \ell\mid m.
\]
Write
\[
    m=\ell r.
\]
Then
\[
    x^m-1
    =
    x^{\ell r}-1
    =
    (x^\ell)^r-1,
\]
and hence
\[
    x^m-1
    =
    (x^\ell-1)
    \left(
        1+x^\ell+x^{2\ell}+\cdots+x^{(r-1)\ell}
    \right).
\]
Therefore
\[
    x^\ell-1\divides x^m-1.
\]

Conversely, suppose
\[
    x^\ell-1\divides x^m-1.
\]
By the division algorithm for integers, write
\[
    m=\ell q+r,
    \qquad
    0\le r<\ell.
\]
Then
\begin{align*}
    x^m-1
    &=
    x^{\ell q+r}-x^r+x^r-1\\
    &=
    x^r\frac{x^{\ell q}-1}{x^\ell-1}(x^\ell-1)
    +(x^r-1).
\end{align*}
The first term is divisible by $x^\ell-1$. Since the whole polynomial
$x^m-1$ is assumed divisible by $x^\ell-1$, it follows that
\[
    x^\ell-1\divides x^r-1.
\]
But
\[
    \deg(x^r-1)<\deg(x^\ell-1)
\]
when $0<r<\ell$. Hence this is possible only if
\[
    r=0.
\]
Therefore
\[
    \ell\mid m.
\]
\end{proof}

The same argument with $x$ replaced by an integer $a>1$ shows that, for
positive integers $\ell,m$,
\[
    a^\ell-1\mid a^m-1
    \quad\Longleftrightarrow\quad
    \ell\mid m.
\]

% ============================================================
\section{Subfields of a Finite Field}
% ============================================================

Let $F$ be a finite field of characteristic $p$. Let
\[
    \Z/p\Z\subseteq K\subseteq F.
\]
Write
\[
    |F|=p^n,
    \qquad
    |K|=p^d.
\]

Since $F$ is a vector space over $K$, say of finite dimension $r$,
\[
    F\cong K^r
\]
as a $K$-vector space. Hence
\[
    |F|=|K|^r.
\]
Therefore
\[
    p^n=(p^d)^r=p^{dr},
\]
so
\[
    n=dr.
\]
Consequently,
\[
    \boxed{d\mid n.}
\]

Thus a necessary condition for a subfield of order $p^d$ to exist inside
a field of order $p^n$ is
\[
    d\mid n.
\]

\begin{remark}
For example, if
\[
    |F|=p^3,
\]
then there is no subfield strictly between
\[
    \Z/p\Z
\]
and $F$, since the only positive divisors of $3$ are $1$ and $3$.
\end{remark}

% ============================================================
\section{Three Key Questions}
% ============================================================

The lecture isolates three key questions:

\begin{enumerate}[label=\arabic*.]
    \item Can we always find a subfield of order $p^d$ given
    \[
        d\mid n?
    \]

    \item How many such fields can we get?

    \item For a given prime $p$ and integer $n>1$, does a finite field
    of order
    \[
        p^n
    \]
    exist?
\end{enumerate}

% ============================================================
\section{The Subfield Theorem}
% ============================================================

\begin{theorem}
Let $F$ be a finite field of order
\[
    p^n.
\]
Then the subfields of $F$ are in one-to-one correspondence with the
divisors of $n$.

In particular, for every
\[
    d\mid n,
\]
there is a unique subfield of $F$ having
\[
    p^d
\]
elements.
\end{theorem}

\subsection{Construction of the Subfield}

Let
\[
    d\mid n.
\]
Define
\[
    K
    =
    \{\alpha\in F:\alpha^{p^d}=\alpha\}.
\]
Equivalently, $K$ is the set of roots in $F$ of
\[
    x^{p^d}-x.
\]

We first show that $K$ is a field.

Let
\[
    \alpha,\beta\in K.
\]
Then
\[
    \alpha^{p^d}=\alpha,
    \qquad
    \beta^{p^d}=\beta.
\]
Using Freshman's dream,
\[
    (\alpha+\beta)^{p^d}
    =
    \alpha^{p^d}+\beta^{p^d}
    =
    \alpha+\beta.
\]
Hence
\[
    \alpha+\beta\in K.
\]
In particular $0=\alpha+(-\alpha)$ will lie in $K$ once we check additive
inverses. Freshman's dream gives
\[
    0=(\alpha+(-\alpha))^{p^d}
    =
    \alpha^{p^d}+(-\alpha)^{p^d},
\]
so
\[
    (-\alpha)^{p^d}=-\alpha^{p^d}=-\alpha.
\]
Thus $-\alpha\in K$.

Similarly,
\[
    (\alpha\beta)^{p^d}
    =
    \alpha^{p^d}\beta^{p^d}
    =
    \alpha\beta,
\]
so
\[
    \alpha\beta\in K.
\]

If $\alpha\neq0$, then
\[
    (\alpha^{-1})^{p^d}
    =
    (\alpha^{p^d})^{-1}
    =
    \alpha^{-1},
\]
so
\[
    \alpha^{-1}\in K.
\]
Therefore $K$ is a field.

% ============================================================
\section{Why $K$ Has $p^d$ Elements}
% ============================================================

All elements of $K$ are roots of
\[
    x^{p^d}-x.
\]
Conversely, every root of this polynomial in $F$ belongs to $K$.

Since
\[
    x^{p^d}-x
\]
has degree $p^d$, it has at most $p^d$ roots. The lecture shows that,
for $d\mid n$, all its roots occur in $F$, so
\[
    |K|=p^d.
\]

Indeed, since
\[
    d\mid n,
\]
we have
\[
    p^d-1\mid p^n-1.
\]
Using the divisibility result,
\[
    x^{p^d}-x
    \divides
    x^{p^n}-x.
\]
The polynomial $x^{p^n}-x$ has derivative $-1$ in characteristic $p$, so it
is square-free and splits into $p^n$ distinct linear factors in $F$. Any
divisor is therefore also square-free and splits completely in $F$. Hence
the roots of
\[
    x^{p^d}-x
\]
lie in $F$.

Therefore
\[
    \boxed{|K|=p^d.}
\]

% ============================================================
\section{Uniqueness of the Subfield}
% ============================================================

Let $K'$ be any other \emph{subfield of $F$} with $p^d$ elements.

Since $K'$ is a finite field of characteristic $p$,
every $\alpha\in K'$ satisfies
\[
    \alpha^{p^d}-\alpha=0.
\]
Hence every element of $K'$ is a root of
\[
    x^{p^d}-x.
\]
But $K$ consists precisely of the roots of this polynomial in $F$.
Therefore
\[
    K'\subseteq K,
\]
and comparing cardinalities gives $K'=K$.

Thus, inside a fixed finite field $F$, the subfield of order $p^d$ is
unique.

\begin{corollary}
The subfields of a finite field $F$ of order $p^n$ correspond exactly to
the divisors of $n$:
\[
    d\mid n
    \quad\longleftrightarrow\quad
    \text{unique subfield of order }p^d.
\]
\end{corollary}

% ============================================================
\section{Constructing a Finite Field of Order $p^n$}
% ============================================================

\begin{theorem}
Let $K$ be a field and let
\[
    f(x)\in K[x]
\]
be an irreducible polynomial of degree $n$. Then there exists a field
$F$ containing $K$ and a zero $\alpha$ of $f(x)$ such that
\[
    [F:K]=n.
\]
\end{theorem}

\begin{proof}
Since $K[x]$ is a Euclidean domain, it is a PID and hence a UFD.
Because $f(x)$ is irreducible,
\[
    (f(x))
\]
is a maximal ideal of $K[x]$. Therefore
\[
    \boxed{K[x]/(f(x))}
\]
is a field.

Let
\[
    F=K[x]/(f(x)).
\]
The natural embedding
\[
    K\hookrightarrow F
\]
identifies $K$ with a subfield of $F$.

Let
\[
    \alpha=x+(f(x))
\]
be the residue class of $x$. Then
\[
    f(\alpha)=0
\]
in $F$.

Every element of $F$ has a unique representative of degree less than $n$:
\[
    a_0+a_1\alpha+\cdots+a_{n-1}\alpha^{n-1},
    \qquad
    a_i\in K.
\]
Hence
\[
    1,\alpha,\ldots,\alpha^{n-1}
\]
form a basis of $F$ over $K$, and therefore
\[
    [F:K]=n.
\]
\end{proof}

% ============================================================
\section{Irreducible Polynomials and Field Construction}
% ============================================================

The goal is therefore:

\begin{center}
\boxed{\text{Construct a field with }p^n\text{ elements.}}
\end{center}

Start with
\[
    \Fp=\Z/p\Z
\]
and find an irreducible polynomial of degree $n$:
\[
    f(x)\in\Fp[x].
\]
Then
\[
    \boxed{
    \Fp[x]/(f(x))
    }
\]
is a field with
\[
    p^n
\]
elements.

Indeed, each element has a unique representation
\[
    a_0+a_1x+\cdots+a_{n-1}x^{n-1}+(f(x)),
    \qquad
    a_i\in\Fp,
\]
so there are
\[
    p^n
\]
such elements.

\subsection{Examples}

Over $\Z/2\Z$,
\[
    x^2+x+1
\]
is irreducible in
\[
    (\Z/2\Z)[x].
\]
Therefore
\[
    \boxed{
    \frac{(\Z/2\Z)[x]}{(x^2+x+1)}
    }
\]
is a field with
\[
    2^2=4
\]
elements.

Over $\Z/3\Z$,
\[
    x^2+1
\]
is irreducible. Hence
\[
    \boxed{
    \frac{(\Z/3\Z)[x]}{(x^2+1)}
    }
\]
is a field with
\[
    3^2=9
\]
elements.

% ============================================================
\section{Counting Irreducible Polynomials}
% ============================================================

The next aim is to count irreducible polynomials.

Let
\[
    N_d
\]
denote the number of monic irreducible polynomials in
\[
    \Fp[x]
\]
of degree $d$.

The lecture uses the factorization of
\[
    x^{p^n}-x
\]
to count these polynomials.

\begin{theorem}
For every $n\ge1$,
\[
    \boxed{
    x^{p^n}-x
    =
    \prod_{d\mid n}
    \prod_{\deg f=d} f(x),
    }
\]
where the inner product is over all monic irreducible polynomials
$f(x)\in\Fp[x]$ of degree $d$.
\end{theorem}

\subsection{Why the Irreducible Factors Have Degree Dividing $n$}

Let
\[
    f(x)
\]
be an irreducible polynomial of degree $d$ dividing
\[
    x^{p^n}-x.
\]
Let
\[
    \alpha
\]
be a zero of $f(x)$, and let
\[
    K=\Fp[x]/(f(x)).
\]
Then
\[
    |K|=p^d.
\]

Since $f(x)\mid x^{p^n}-x$, we have
\[
    \alpha^{p^n}=\alpha
\]
in $K$. The Frobenius map $\beta\mapsto\beta^{p^n}$ is a field automorphism
of $K$ (by Freshman's dream), and it fixes $\alpha$. Because
$K=\Fp(\alpha)$, the map is the identity on $K$. Thus every $\beta\in K$
satisfies $\beta^{p^n}=\beta$. Every element of $K$ is therefore a root of
the square-free polynomial $x^{p^d}-x$, and also a root of $x^{p^n}-x$.
Hence
\[
    x^{p^d}-x
    \divides
    x^{p^n}-x
\]
in $\Fp[x]$. Therefore
\[
    p^d-1\mid p^n-1,
\]
and the integer divisibility criterion already proved yields
\[
    d\mid n.
\]

Therefore every irreducible factor of
\[
    x^{p^n}-x
\]
has degree dividing $n$.

\subsection{Conversely}

Suppose
\[
    d\mid n
\]
and let
\[
    f(x)
\]
be an irreducible polynomial of degree $d$ over $\Fp$.

Let $\alpha$ be a root of $f(x)$. If $\alpha=0$, then $f(x)=x$ (up to units),
and $x$ divides $x^{p^n}-x$. Now assume $\alpha\neq0$. Then
\[
    \Fp(\alpha)
\]
has
\[
    p^d
\]
elements.

Since
\[
    d\mid n,
\]
we have
\[
    p^d-1\mid p^n-1.
\]
Thus
\[
    \alpha^{p^n-1}=1,
\]
and hence
\[
    \alpha^{p^n}=\alpha.
\]
Therefore
\[
    f(x)\mid x^{p^n}-x.
\]

Consequently,
\[
    \boxed{
    f(x)\mid x^{p^n}-x
    \quad\Longleftrightarrow\quad
    \deg f\mid n
    }
\]
for irreducible $f(x)\in\Fp[x]$.

% ============================================================
\section{Counting Formula}
% ============================================================

Since
\[
    x^{p^n}-x
\]
is the product of all monic irreducible polynomials whose degrees divide
$n$, taking degrees gives
\[
    p^n
    =
    \sum_{d\mid n} dN_d.
\]

Thus
\[
    \boxed{
    p^n=\sum_{d\mid n}dN_d.
    }
\]

This is the fundamental counting identity for irreducible polynomials
over $\Fp$.

% ============================================================
\section{Möbius Inversion}
% ============================================================

For $n>0$, Möbius inversion gives
\[
    N_n
    =
    \frac1n
    \sum_{d\mid n}
    \mu(d)p^{n/d}.
\]

Equivalently,
\[
    \boxed{
    N_n
    =
    \frac1n
    \sum_{d\mid n}
    \mu(d)p^{n/d}.
    }
\]

Writing the terms explicitly,
\[
    N_n
    =
    \frac1n
    \left(
        p^n
        +\sum_{\substack{d\mid n\\d>1}}
        \mu(d)p^{n/d}
    \right).
\]

The lecture then observes that the correction term is strictly smaller
than $p^n$, which proves positivity. In full:
\[
    N_n
    \ge
    \frac1n
    \Biggl(
        p^n
        -
        \sum_{\substack{d\mid n\\ d>1}}
        p^{n/d}
    \Biggr).
\]
For $n=1$ one has $N_1=p>0$. For $n>1$ each proper exponent $n/d$ is at
most $n-1$, so
\[
    \sum_{\substack{d\mid n\\ d>1}}
    p^{n/d}
    \le
    \sum_{j=1}^{n-1}p^j
    =
    p\frac{p^{n-1}-1}{p-1}
    <p^n.
\]
Thus $N_n>0$.

Hence for every prime $p$ and every $n\ge1$, there exists an irreducible
polynomial of degree $n$ over $\Fp$.

% ============================================================
\section{A Useful Bound}
% ============================================================

From
\[
    N_n
    =
    \frac1n
    \sum_{d\mid n}\mu(d)p^{n/d},
\]
we obtain
\[
    nN_n
    \le
    p^n.
\]

The identity $p^n=\sum_{d\mid n}dN_d$ together with $N_n>0$ (already
proved) immediately gives, for $n>1$,
\[
    \sum_{\substack{d\mid n\\d<n}}dN_d
    =
    p^n-nN_n
    <p^n.
\]
The lecture also sketches a direct estimate not using M\"obius inversion.

It also records the elementary estimate
\[
    N_d\le p^d.
\]
Hence
\[
    \sum_{\substack{d\mid n\\d<n}}dN_d
    \le
    \sum_{\substack{d\mid n\\d<n}}dp^d.
\]

Since every proper divisor satisfies
\[
    d\le\frac n2,
\]
the lecture bounds the sum by a geometric-series type expression and
obtains an upper bound of the form
\[
    \sum_{\substack{d\mid n\\d<n}}dp^d
    <
    2n\,p^{n/2}.
\]
This last inequality may be checked directly for $2\le n\le 9$, while for
$n\ge 10$ one has $2n<p^{n/2}$ and therefore
\[
    \sum_{\substack{d\mid n\\d<n}}dp^d
    <p^n.
\]

The purpose of this estimate is to show directly that the leading term
$p^n$ in Möbius inversion dominates all the lower-order contributions.

% ============================================================
\section{Factorization of $x^{p^n}-x$}
% ============================================================

The polynomial
\[
    x^{p^n}-x
\]
is a product of irreducible polynomials.

Let
\[
    f(x)
\]
be an irreducible factor. Since
\[
    x^{p^n}-x
\]
has derivative
\[
    \frac{d}{dx}(x^{p^n}-x)
    =
    p^n x^{p^n-1}-1
    =
    -1
\]
in characteristic $p$, it has no repeated roots.

Equivalently, if
\[
    x^{p^n}-x=g(x)^2h(x),
\]
then differentiation would give
\[
    p^n x^{p^n-1}-1
    =
    2g(x)g'(x)h(x)+g(x)^2h'(x),
\]
and the right-hand side would be divisible by $g(x)$, whereas the left
side is
\[
    -1.
\]
This is impossible.

Thus
\[
    x^{p^n}-x
\]
is square-free.

% ============================================================
\section{Final Characterization of the Irreducible Factors}
% ============================================================

Let
\[
    f(x)
\]
be an irreducible polynomial of degree $d$ over $\Fp$.

Then
\[
    f(x)\mid x^{p^n}-x
\]
if and only if
\[
    d\mid n.
\]

\begin{proof}
Suppose first that
\[
    f(x)\mid x^{p^n}-x.
\]
Let $\alpha$ be a root of $f(x)$. Then
\[
    \alpha^{p^n}=\alpha.
\]
Write $K=\Fp[x]/(f(x))\cong\Fp(\alpha)$, so $|K|=p^d$. The Frobenius map
$\beta\mapsto\beta^{p^n}$ is a field automorphism of $K$ fixing $\alpha$,
hence the identity on $K=\Fp(\alpha)$. Thus every element of $K$ satisfies
$X^{p^n}-X=0$, so $x^{p^d}-x$ divides $x^{p^n}-x$. Therefore
\[
    p^d-1\mid p^n-1,
\]
and we conclude
\[
    d\mid n.
\]

Conversely, suppose
\[
    d\mid n.
\]
Since
\[
    \Fp(\alpha)
\]
has $p^d$ elements,
\[
    \alpha^{p^d}=\alpha.
\]
Iterating this identity,
\[
    \alpha^{p^n}=\alpha
\]
because $d\mid n$. Therefore $\alpha$ is a root of
\[
    x^{p^n}-x.
\]
Since $f(x)$ is the minimal polynomial of $\alpha$ over $\Fp$,
\[
    f(x)\mid x^{p^n}-x.
\]
\end{proof}

Thus
\[
    \boxed{
    x^{p^n}-x
    =
    \prod_{d\mid n}
    \prod_{\substack{f\ \mathrm{monic}\\
                      f\ \mathrm{irreducible}\\
                      \deg f=d}}
    f(x).
    }
\]

% ============================================================
\section*{Summary of the Lecture}
% ============================================================

The lecture develops the following chain:
\[
\boxed{
\text{Finite fields}
\longrightarrow
\text{Freshman's dream}
\longrightarrow
\text{subfields}
\longrightarrow
\text{irreducible polynomials}
\longrightarrow
\text{counting irreducibles}
}.
\]

The main results are:

\[
    (\alpha+\beta)^{p^n}
    =
    \alpha^{p^n}+\beta^{p^n},
\]
\[
    x^\ell-1\mid x^m-1
    \quad\Longleftrightarrow\quad
    \ell\mid m,
\]
\[
    |F|=p^n,\quad K\subseteq F
    \quad\Longrightarrow\quad
    |K|=p^d\text{ for some }d\mid n,
\]
\[
    d\mid n
    \quad\Longleftrightarrow\quad
    \text{there is a unique subfield of order }p^d,
\]
\[
    \Fp[x]/(f(x))
    \text{ is a field whenever }f(x)\text{ is irreducible},
\]
\[
    |\Fp[x]/(f(x))|=p^{\deg f},
\]
\[
    p^n=\sum_{d\mid n}dN_d,
\]
and
\[
    \boxed{
    N_n
    =
    \frac1n\sum_{d\mid n}\mu(d)p^{n/d}.
    }
\]

Finally,
\[
    \boxed{
    f(x)\mid x^{p^n}-x
    \quad\Longleftrightarrow\quad
    \deg f\mid n
    }
\]
for every irreducible polynomial $f(x)\in\Fp[x]$.

\end{document}
